\documentclass[11pt]{article}
\usepackage{tcee}
\addbibresource{3B28.bib}
\tcetema[Economía de los países en desarrollo y teorías del desarrollo]{3.B.28}{Economía de los países en desarrollo. Teorías recientes del desarrollo económico. Evidencia empírica, con especial referencia a la aproximación experimental, e implicaciones para el diseño de políticas.}
\tcefecha{11/11/2026}
\begin{document}
\tceetitulo

\section*{Introducción}
\minutos{2,5}

\begin{itemize}
\item \textbf{\emph{Enganche}}:
  \begin{itemize}
  \item «Unos nacen con estrella y otros estrellados», dice el refrán. En la economía mundial, la estrella tiene nombre: el país en el que se nace.
  \item Un niño que nace hoy en Kenia puede esperar vivir unos 64 años; uno que nace en España, unos 84: 20 años más. Y la renta por habitante de España, en paridad de poder adquisitivo, es unas 8 veces la de Kenia \parencite{PNUD2025IDH}.
  \item Ninguno de los 2 eligió dónde nacer. Para \textcite{Milanovic2016}, la mayor parte de la desigualdad de renta entre los habitantes del mundo se debe al país en el que viven, una \emph{prima de ciudadanía} que nadie se gana.
  \item Detrás está la pregunta de esta exposición: ¿por qué unos países siguen siendo pobres, y qué sabemos de verdad sobre cómo salir de la pobreza?
  \end{itemize}

\item \textbf{\emph{Relevancia}}:
  \begin{itemize}
  \item \emph{La pobreza}: en 2024, el 10,1~\% de la población mundial, unos 824 millones de personas, vivía con menos de 3 dólares al día (en paridad de poder adquisitivo de 2021), frente al 43,4~\% en 1990 \parencite{BM2026PobrezaSept,BM2026PIP}. La caída es el mayor logro económico de nuestro tiempo, pero 7 de cada 10 pobres extremos viven hoy en el África subsahariana.
  \item \emph{La convergencia de una parte del mundo}: el PIB real por habitante de las economías emergentes y en desarrollo se ha multiplicado por 2,6 entre 1995 y 2024, frente a 1,5 en las avanzadas \parencite{FMI2026WEO}. Unos países han salido del atraso y otros no: entender por qué es la tarea de la economía del desarrollo.
  \item \emph{El reconocimiento}: la disciplina ha dado una lista larga de premios Nobel: \autor{Myrdal} (1974), \autor{Lewis} y \autor{Schultz} (1979), \autor{North} (1993), \autor{Sen} (1998), \autor{Deaton} (2015), \autor{Banerjee}, \autor{Duflo} y \autor{Kremer} (2019), «por su enfoque experimental para aliviar la pobreza global» \parencite[traducción propia]{NobelPrize2019}, y \autor{Acemoglu}, \autor{Johnson} y \autor{Robinson} (2024), por sus estudios sobre cómo se forman las instituciones y cómo afectan a la prosperidad \parencite{NobelPrize2024}. Los 2 últimos miran al desarrollo desde los 2 extremos: el programa concreto y las reglas del juego.
  \end{itemize}

\item \textbf{\emph{Contextualización}}:
  \begin{itemize}
  \item Desde un punto de vista conceptual:
  \begin{itemize}
  \item El \concepto{crecimiento económico} es el aumento sostenido de la renta por habitante: una dimensión \emph{cuantitativa} \vertema{3.A.43}.
  \item El \concepto{desarrollo económico} es la transformación de la estructura productiva, social e institucional que amplía lo que las personas pueden ser y hacer, sus \emph{capacidades} en el sentido de \textcite{Sen1999}: una dimensión \emph{cualitativa}.
  \item El crecimiento suele ser condición \emph{necesaria} del desarrollo, porque da los recursos, pero no \emph{suficiente}. Guinea Ecuatorial y Sri Lanka tienen casi la misma renta por habitante (unos 12.700 dólares de paridad de poder adquisitivo de 2021), pero Sri Lanka ocupa el puesto 89 en desarrollo humano y Guinea Ecuatorial, el 133 \parencite{PNUD2025IDH}.
  \item Un \concepto{país en desarrollo} es, a efectos prácticos, uno de renta baja o media en la clasificación del Banco Mundial; en ellos vive más del 80~\% de la población mundial \parencite{BM2026WDI}.
  \end{itemize}
  \item Desde un punto de vista histórico:
  \begin{itemize}
  \item La pregunta por las causas de la riqueza de las naciones es la que fundó la economía \parencite{Smith1776Wealth}, y su marco moderno es la teoría del crecimiento, desde \textcite{Solow1956}.
  \item A la vez que \autor{Solow}, tras la Segunda Guerra Mundial y con la descolonización, nace la \concepto{economía del desarrollo} como rama propia: sus fundadores pensaban que los países pobres tenían peculiaridades (dualidad, especialización en productos primarios, instituciones débiles, mercados incompletos) que pedían un enfoque distinto.
  \item Desde entonces ha pasado por 3 etapas: el optimismo intervencionista de los 50 y 60, en el que el Estado debía empujar la industrialización; la contrarrevolución neoclásica de los 80, con el mercado y la estabilidad macroeconómica; y, desde los 90, unas teorías que formalizan las viejas intuiciones (equilibrios múltiples, instituciones) y una revolución empírica, la de los experimentos.
  \item Por eso no hay un programa de investigación asentado, como en la teoría del consumo: conviven los \emph{modelos estructurales}, que van de la teoría a los datos, y el trabajo \emph{experimental}, que va de los datos a la teoría, sin acuerdo sobre cuál debe guiar la política \vertema{3.A.1}.
  \end{itemize}
  \end{itemize}

\item \textbf{\emph{Problemática (preguntas clave)}}:
  \begin{itemize}
  \item ¿Qué tienen en común los países en desarrollo y cómo se mide su atraso?
  \item ¿Por qué hacía falta una teoría del desarrollo distinta de la del crecimiento, y qué queda de ella?
  \item ¿Qué explican las teorías recientes y qué dice la evidencia, en especial la experimental, sobre lo que funciona contra la pobreza?
  \item ¿Qué implica todo ello para el diseño de las políticas?
  \end{itemize}

\item \textbf{\emph{Estructura}}:
  \begin{itemize}
  \item Tras fijar los hechos, la exposición mira el desarrollo desde 2 escalas: la del país, que pregunta por qué unas economías se quedan atrás, y la de las personas, que pregunta cómo viven, deciden y salen de la pobreza los hogares pobres, y qué programas lo consiguen.
  \begin{itemize}
  \item \concepto{I. Hechos estilizados} (\ref{sec:desarrollo}): cómo se mide el desarrollo y qué rasgos comparten los países en desarrollo.
  \item \concepto{II. Aproximación macroeconómica} (\ref{sec:macro}): las tendencias del crecimiento, la desigualdad y la pobreza, y las teorías, de las pioneras a las recientes, sobre por qué algunos países se quedan atrás.
  \item \concepto{III. Aproximación microeconómica} (\ref{sec:microap}): de \autor{Sen} y \autor{Deaton} a los mercados de los pobres y la aproximación experimental, y las implicaciones para el diseño de políticas.
  \end{itemize}
  \end{itemize}
\end{itemize}

\begin{esquema}
\textbf{3.B.28. Economía de los países en desarrollo, teorías recientes del desarrollo, evidencia experimental y políticas}\\[2pt]
Introducción\esqtiempo{2,5}\\
I. Hechos estilizados de las economías en desarrollo\esqtiempo{3}\\
\hspace*{1.5em}1. Medición del desarrollo\\
\hspace*{1.5em}2. Características de los países en desarrollo\\
II. Aproximación macroeconómica\esqtiempo{10}\\
\hspace*{1.5em}1. Tendencias en crecimiento, desigualdad y pobreza\\
\hspace*{1.5em}2. Las teorías pioneras y su replanteamiento\\
\hspace*{1.5em}3. Teorías recientes: complementariedades e instituciones\\
III. Aproximación microeconómica\esqtiempo{13,5}\\
\hspace*{1.5em}1. Las contribuciones pioneras: \autor{Sen} y \autor{Deaton}\\
\hspace*{1.5em}2. Comportamiento y mercados de los pobres\\
\hspace*{1.5em}3. La aproximación experimental\\
\hspace*{1.5em}4. Implicaciones para el diseño de políticas\\
Conclusión\esqtiempo{1}
\end{esquema}

\begin{notaopositor}
Reparto del cante: introducción, \tiempo{2,5}; I, \tiempo{3} (I.1, \tiempo{1,5}; I.2, \tiempo{1,5}); II, \tiempo{10} (II.1, \tiempo{1,5}; II.2, \tiempo{5,5}; II.3, \tiempo{3}); III, \tiempo{13,5} (III.1, \tiempo{2}; III.2, \tiempo{3}; III.3, \tiempo{6}; III.4, \tiempo{2,5}), y conclusión, \tiempo{1}. El título pide sobre todo lo reciente: a las teorías desde los 90 (apartados~\ref{sec:recientes} y~\ref{sec:microap}), la aproximación experimental y las políticas hay que dedicarles al menos un cuarto de hora. Los gráficos que se cantan son el mercado de trabajo de \autor{Lewis} y la trampa de pobreza (apartado~\ref{sec:clasicas}), y la función clave, la descomposición del sesgo de selección \eqref{eq:sesgo-seleccion} (apartado~\ref{sec:experimental}). Las pegatinas: «una golondrina no hace verano» (el gran empujón), «café para todos» (el Consenso de Washington), «el hábito no hace al monje» (las instituciones de segundo óptimo), «ver para creer» (los experimentos), «del dicho al hecho hay mucho trecho» (la validez externa) y «obras son amores, que no buenas razones» (las políticas). Lo que va en cajas grises, en notas y en los anexos es para las preguntas del tribunal.
\end{notaopositor}

\begin{anotaciones}
\textbf{Notación.} $Y$, $K$, $L$: producción, capital y trabajo; lo que es \emph{por habitante} se escribe como cociente, $Y/L$, y su tasa de crecimiento, $\gamma_{Y/L}$, con $\gamma_X\equiv\dot X/X$. Como en el resto del temario, $\tilde k$ e $\tilde y$ (con virgulilla) son magnitudes por unidad de trabajo eficiente y $k$ e $y$ (sin ella), por habitante \vertema{3.A.43}. En el modelo de \autor{Lewis}, $L_M$ es el empleo del sector moderno, $\overline w_A$ el salario institucional del sector tradicional y $w_M$ el salario del moderno; $w$ es siempre un salario. $v\equiv K/Y$ es la relación capital-producto. Las siglas van en redonda: RNB (renta nacional bruta), PPA (paridad de poder adquisitivo), IDH (índice de desarrollo humano), RRI (relación real de intercambio) e IG (índice de \autor{Gini}). En el apartado~\ref{sec:experimental}, $Y_i(1)$ e $Y_i(0)$ son los resultados potenciales de la persona $i$ con y sin el programa, y $D_i$ indica si lo recibe: ahí $Y$ no es la renta. Las demás letras son locales y se dicen donde salen.
\end{anotaciones}

\begin{anotaciones}
\textbf{Para saber más}
\begin{itemize}
\item \autor{Abhijit Banerjee} y \autor{Esther Duflo}, \emph{Poor Economics}: la economía de la pobreza contada desde los experimentos, con ejemplos de Kenia, la India o Marruecos \parencite{BanerjeeDuflo2011}.
\item \autor{Debraj Ray}, \emph{Development Economics}: el manual de referencia de la disciplina, con los modelos de \autor{Lewis}, de las trampas y de los mercados incompletos \parencite{Ray1998}.
\item \autor{Filipe Campante}, \autor{Federico Sturzenegger} y \autor{Andrés Velasco}, \emph{Advanced Macroeconomics: An Easy Guide}, libro de acceso abierto; su capítulo sobre las trampas de pobreza resume la evidencia: \url{https://doi.org/10.31389/lsepress.ame} \parencite{Campante2021}.
\end{itemize}
\end{anotaciones}

\section{Hechos estilizados de las economías en desarrollo}\label{sec:desarrollo}
\minutos{3}

\begin{notaopositor}
En el cante (\tiempo{3}): la medición (\tiempo{1,5}), con el hilo de las 3 medidas (el desarrollo humano como medida de referencia, la renta como su componente y base de la clasificación, la pobreza como su cara más dura) y las cifras justas: 3 dimensiones del IDH, 4 grupos de renta, la línea de 3 dólares; las características (\tiempo{1,5}), con los 3 grupos de rasgos. Los mapas (gráficos~\ref{fig:mapa-idh} y~\ref{fig:mapa-renta}) y la tabla~\ref{tab:hechos} son de apoyo.
\end{notaopositor}

\begin{itemize}
\item \emph{Conclusión anticipada}: los países en desarrollo son muy distintos entre sí, pero comparten unos rasgos que se pueden medir y que se refuerzan unos a otros. Como el desarrollo tiene varias dimensiones, ninguna cifra lo resume: se mide con un índice compuesto, el de desarrollo humano, y se completa con la renta y con la pobreza.
\item \emph{Nexo}: antes de las teorías, los hechos. Empezamos por cómo se mide el desarrollo.
\end{itemize}

% --------------------------------------------------------------------------
\subsection{Medición del desarrollo}\label{sec:medicion}

\begin{itemize}
\item \emph{El punto de partida}: si el desarrollo es ampliar lo que las personas pueden ser y hacer (introducción), hay que medir a la vez cuánto tienen, cómo viven y quién se queda atrás. De ahí 3 medidas que se complementan: el \emph{desarrollo humano}, que las resume; la \emph{renta}, que es su base material y la que usa la clasificación oficial de los países, y la \emph{pobreza}, que mira a los que están peor.
\end{itemize}

\paragraph{El desarrollo humano: la medida de referencia}

\begin{itemize}
\item \emph{El índice}: el \concepto{índice de desarrollo humano} (IDH), que publica el Programa de las Naciones Unidas para el Desarrollo (PNUD) desde 1990, es el indicador de referencia del desarrollo. Nace del enfoque de las capacidades de \autor{Sen}: la renta es solo un medio \parencite{PNUD1990Human,Sen1985Commodities} \vertema{3.A.24}. Combina 3 dimensiones:
  \begin{itemize}
  \item \emph{salud}: la esperanza de vida al nacer;
  \item \emph{educación}: los años esperados de escolarización de un niño y los años medios de escolarización de los adultos;
  \item \emph{nivel de vida}: la RNB por habitante en PPA, en logaritmos, porque cada dólar adicional amplía menos las capacidades.
  \end{itemize}
\item \emph{La agregación}: cada dimensión se normaliza entre 0 y 1 y el IDH es su \emph{media geométrica}, no la aritmética: así, una dimensión muy baja no se compensa del todo con las otras \parencite{PNUD2025NotasTecnicas}. El cálculo, paso a paso, está en el anexo~\ref{anexo:idh}.
\item \emph{Los 4 niveles}: el PNUD clasifica los países en desarrollo humano muy alto (IDH de 0,800 o más), alto, medio y bajo (menos de 0,550). En 2023 lo encabeza Islandia (0,972) y lo cierra Sudán del Sur (0,388); España ocupa el puesto 28 (0,918). El mapa (gráfico~\ref{fig:mapa-idh}) muestra la geografía del atraso: casi todos los países de desarrollo humano bajo están en el África subsahariana \parencite{PNUD2025IDH}.
\end{itemize}

\begin{figure}[H]
\centering
\caption{Niveles de desarrollo humano del PNUD, 2023}\label{fig:mapa-idh}
\grafico[0.95\textwidth]{mapa-idh}
\fuente{elaboración propia a partir de \textcite{PNUD2025Series} (Informe sobre Desarrollo Humano 2025, IDH de 2023). Fronteras de Natural Earth (1:110 millones); mapa generado con \texttt{graficos/mapas-desarrollo.py}.}
\end{figure}

\paragraph{La renta: la base material y la clasificación de los países}

\begin{itemize}
\item \emph{El indicador}: la \concepto{renta nacional bruta} (RNB) por habitante, que suma al PIB las rentas netas recibidas del exterior, es la dimensión material del IDH y el criterio con el que se clasifica a los países. Se prefiere la renta \emph{nacional} a la producción \emph{interior} porque un país con mucha inversión extranjera puede producir mucho y quedarse con poco: las filiales repatrían sus beneficios.
\item \emph{La clasificación del Banco Mundial}: cada 1 de julio ordena los países por su RNB por habitante del año anterior en 4 grupos: renta baja, media-baja, media-alta y alta (gráfico~\ref{fig:mapa-renta}). Desde julio de 2026, la renta baja llega hasta 1.175 dólares y la alta empieza en 14.375 \parencite{BM2026ClasificacionRenta}; los países de renta baja y media son los países en desarrollo de la introducción.
\item \emph{Por qué no basta}: la renta es un medio y una media. No dice cómo se reparte ni en qué se convierte: con casi la misma renta por habitante, un país puede tener un desarrollo humano mucho mayor que otro (introducción). Por eso el IDH la combina con la salud y la educación, y la pobreza mira a quienes quedan abajo.
\end{itemize}

\begin{figure}[H]
\centering
\caption{Grupos de renta del Banco Mundial, desde julio de 2026}\label{fig:mapa-renta}
\grafico[0.95\textwidth]{mapa-renta}
\fuente{elaboración propia a partir de \textcite{BM2026ClasificacionRenta} (clasificación del ejercicio 2027, con la RNB por habitante de 2025). Fronteras de Natural Earth (1:110 millones); mapa generado con \texttt{graficos/mapas-desarrollo.py}.}
\end{figure}

\begin{anotaciones}
\textbf{La renta en detalle.}
\begin{itemize}
\item \emph{El método Atlas y la PPA}: la clasificación usa tipos de cambio de mercado suavizados (método Atlas). A tipos de cambio de mercado se infravalora la renta de los países pobres, donde los servicios que no se comercian son más baratos; por eso, para comparar niveles de vida, se usa la \concepto{paridad de poder adquisitivo} (PPA): la renta por habitante de los países de renta alta es 65 veces la de los de renta baja a tipos de cambio de mercado y 26 veces en PPA \parencite{BM2026WDI}.
\item \emph{Los umbrales de julio de 2026}: renta baja, hasta 1.175 dólares (25 economías); media-baja, de 1.176 a 4.635 (47); media-alta, de 4.636 a 14.375 (59); alta, más de 14.375 (86) \parencite{BM2026ClasificacionRenta}.
\end{itemize}
\end{anotaciones}

\paragraph{La pobreza: la cara más dura del subdesarrollo}

\begin{itemize}
\item \emph{Absoluta y relativa}: la \concepto{pobreza absoluta} es no alcanzar un mínimo fijo de consumo; la \concepto{relativa}, quedar lejos de la mediana del propio país (en la Unión Europea, por debajo del 60~\% de la renta mediana). En los países en desarrollo importa sobre todo la absoluta.
\item \emph{La línea internacional}: desde junio de 2025, el Banco Mundial considera en \concepto{pobreza extrema} a quien vive con menos de 3,00 dólares diarios de PPA de 2021 (antes, 2,15 dólares de PPA de 2017). Con las líneas típicas de los países de renta media-baja y media-alta, 4,20 y 8,30 dólares, es pobre el 18~\% y el 46~\% de la población mundial \parencite{BM2026PIP}.
\item \emph{La pobreza multidimensional}: igual que el IDH completa la renta, el \concepto{índice de pobreza multidimensional}, del PNUD y la Universidad de Oxford, completa la línea de pobreza: cuenta las carencias de cada hogar en salud, educación y condiciones de vida (agua, electricidad, vivienda). Unos 1.100 millones de personas sufren pobreza multidimensional aguda, la mitad de ellas menores \parencite{PNUD2025IPM}.
\item \emph{Nexo}: medido el desarrollo, vemos qué tienen en común los países que no lo han alcanzado. Su evolución en el tiempo, la de la pobreza incluida, es el punto de partida del apartado~\ref{sec:macro}.
\end{itemize}

% --------------------------------------------------------------------------
\subsection{Características de los países en desarrollo}\label{sec:caracteristicas}

\begin{itemize}
\item \emph{Nexo}: los países en desarrollo son muy heterogéneos (Vietnam y Níger no se parecen en casi nada) y no todos tienen todos los rasgos, pero la mayoría comparte 3 grupos de ellos: económicos, institucionales y demográficos. La tabla~\ref{tab:hechos} los resume con datos.
\end{itemize}

\begin{table}[H]
\centering
\caption{Algunos rasgos de los países según su nivel de renta}\label{tab:hechos}
\footnotesize
\renewcommand{\arraystretch}{1.25}
\begin{tabular}{>{\raggedright\arraybackslash}m{6.6cm}>{\centering\arraybackslash}m{2cm}>{\centering\arraybackslash}m{2cm}>{\centering\arraybackslash}m{2cm}}
\toprule
\rowcolor{tceeazulpalido} & \textbf{Renta baja} & \textbf{Renta media-baja} & \textbf{Renta alta}\\
\midrule
RNB por habitante, PPA (dólares internacionales corrientes, 2025) & 2.653 & 10.785 & 69.862\\
Empleo en la agricultura (\% del total, 2025) & 59,4 & 39,8 & 2,9\\
Hijos por mujer (2024) & 4,5 & 2,6 & 1,4\\
Población menor de 15 años (\%, 2025) & 40,9 & 29,5 & 15,5\\
Crecimiento de la población (\% anual, 2025) & 2,7 & 1,3 & 0,4\\
Esperanza de vida al nacer (años, 2024) & 65,2 & 69,5 & 80,4\\
Población urbana (\%, 2025) & 35,2 & 42,0 & 80,2\\
Crédito al sector privado (\% del PIB) & 12,6\textsuperscript{a} & 38,7 & 152,1\textsuperscript{b}\\
Ahorro bruto (\% de la RNB) & 15,8\textsuperscript{b} & 28,2 & 23,2\textsuperscript{b}\\
\bottomrule
\end{tabular}

\smallskip
\fuente{\textcite{BM2026WDI}. Grupos de renta del Banco Mundial. Crédito y ahorro, 2025, salvo (a) 2020 y (b) 2024.}
\end{table}

\paragraph{Rasgos económicos}

\begin{itemize}
\item \emph{Renta baja y dualidad}: un sector tradicional, sobre todo agrario, con mucho empleo y baja productividad, junto a un sector moderno pequeño. En los países de renta baja, 6 de cada 10 ocupados trabajan en la agricultura; en los de renta alta, 3 de cada 100.
\item \emph{Poco paro abierto y mucho subempleo}: sin seguro de desempleo, casi nadie puede permitirse estar parado; el problema es el empleo de baja productividad y la \emph{informalidad}. El 58~\% del empleo mundial es informal, y bastante más en los países de renta baja \parencite{OIT2023Informal}.
\item \emph{Poco ahorro y mercados financieros poco profundos}: el crédito al sector privado es el 13~\% del PIB en los países de renta baja y el 152~\% en los de renta alta, y los seguros apenas existen \parencite{BM2026WDI}.
\item \emph{Dependencia de los productos primarios}: las exportaciones se concentran en materias primas y alimentos, con precios volátiles y una demanda que crece menos que la renta mundial (ley de \autor{Engel}\footnote{La \emph{ley de \autor{Engel}} dice que, al subir la renta, la parte que se gasta en alimentos baja, aunque el gasto en alimentos aumente: su elasticidad renta es menor que 1. Por eso la demanda mundial de productos agrarios crece menos que la renta mundial.}), y las importaciones, en bienes de capital y manufacturas.
  \begin{itemize}
  \item Un auge de las materias primas puede incluso perjudicar: la \concepto{enfermedad holandesa}. Los ingresos del recurso aprecian el tipo de cambio real y atraen capital y trabajo al sector del recurso y a los servicios, y la industria y la agricultura que exportan pierden competitividad \parencite{CordenNeary1982}. Cuando el precio cae, el país se queda sin el recurso y sin la industria.
  \end{itemize}
\item \emph{Vulnerabilidad}: a la deuda externa y a los choques del clima, que golpean más a los países pobres y cálidos \vertema{3.B.29} \vertema{3.B.31}.
\item \emph{Desigualdad}: suele ser más alta que en los países ricos, sobre todo en América Latina y en el África austral; si es una fase del desarrollo o no, es una de las tendencias que se ven en el apartado~\ref{sec:tendencias}.
\end{itemize}

\paragraph{Rasgos institucionales}

\begin{itemize}
\item \emph{Un Estado con poca capacidad}, por exceso o por defecto: recauda poco (en torno al 10~\% del PIB en impuestos en los países de renta baja, según el último dato agregado, de 2016, frente al 15~\% de la administración central en los de renta alta), depende de los aranceles y, a veces, del señoreaje, que alimenta la inflación \parencite{BM2026WDI}.
\item \emph{Reglas del juego débiles}: derechos de propiedad poco seguros, corrupción, mucha economía sumergida y, con frecuencia, regímenes poco democráticos e inestabilidad política, en parte alimentada por la desigualdad. Es el objeto del enfoque institucional (apartado~\ref{sec:institucional}).
\end{itemize}

\paragraph{Rasgos demográficos}

\begin{itemize}
\item \emph{Una transición demográfica a medias}: la mortalidad ya ha caído mucho (la esperanza de vida es de 65 años en los países de renta baja), pero la natalidad sigue alta (4,5 hijos por mujer), de modo que la población crece deprisa (2,7~\% al año) y es muy joven (el 41~\% tiene menos de 15 años) \parencite{BM2026WDI}.
\item \emph{La consecuencia}: hace falta mucha inversión solo para mantener el capital por persona; en el modelo de \autor{Solow}, una $n$ alta baja la renta de estado estacionario \vertema{3.A.43}. No es así en todos: en buena parte de Asia y de América Latina la fecundidad ya está cerca de la de los países ricos.
\end{itemize}

\begin{itemize}
\item \emph{Recapitulación y nexo}: los rasgos de los países en desarrollo se refuerzan entre sí (poca productividad, poco ahorro, poco Estado, mucha población): no son solo síntomas del atraso. La aproximación macroeconómica los convirtió en causas del atraso de los países; la microeconómica, en restricciones de los hogares.
\end{itemize}

% ==========================================================================
\section{Aproximación macroeconómica}\label{sec:macro}
\minutos{10}

\begin{notaopositor}
En el cante (\tiempo{10}): las tendencias (\tiempo{1,5}), de palabra, con 3 cifras de la pobreza (43~\%, 10~\%, 7 de cada 10 en África), la convergencia desde 2000 y la U invertida de \autor{Kuznets}, que los datos no confirman; las teorías pioneras (\tiempo{5,5}): la emergencia y las ideas comunes (\tiempo{0,75}); \autor{Lewis} (\tiempo{1,25}) con el gráfico, que se construye en 3 pasos (gráficos~\ref{fig:lewis-a} a~\ref{fig:lewis}); \autor{Rostow} y la brecha de financiación (\tiempo{0,5}); \autor{Myrdal} y \autor{Kaldor} de palabra (\tiempo{0,5}); \autor{Nurkse}, la trampa y el gran empujón (\tiempo{1,25}), con el gráfico en 3 pasos (gráficos~\ref{fig:trampa-pobreza-a} a~\ref{fig:trampa-pobreza}); la CEPAL y la sustitución de importaciones en resumen (\tiempo{0,5}); el replanteamiento ortodoxo y el Consenso de Washington (\tiempo{0,75}), en orden cronológico. Las teorías recientes (\tiempo{3}): el crecimiento endógeno, los fallos de coordinación, el anillo O y la distancia a la frontera (\tiempo{1,5}), y las instituciones (\tiempo{1,5}). Lo que va en cajas grises y en el anexo~\ref{anexo:clasicas} es de profundización.
\end{notaopositor}

\begin{itemize}
\item \emph{Conclusión anticipada}: a escala de los países, el atraso no se corrige solo. Los primeros teóricos del desarrollo acertaron en las preguntas (cambio estructural, equilibrios múltiples, fallos de coordinación), aunque apenas las formalizaron, y sus recetas, primero intervencionistas y después liberales, fallaron cuando se aplicaron como recetas universales. Las teorías recientes formalizan aquellas intuiciones y apuntan a las instituciones como causa última.
\item \emph{Nexo}: primero, los hechos agregados que hay que explicar; después, las teorías.
\end{itemize}

% --------------------------------------------------------------------------
\subsection{Tendencias globales en crecimiento, desigualdad y pobreza}\label{sec:tendencias}

\paragraph{Crecimiento: una convergencia reciente y parcial}

\begin{itemize}
\item \emph{La convergencia}: no la hubo en el siglo XX, en el que los países ricos se alejaron de los pobres; desde 2000 los pobres crecen algo más deprisa, aunque con fragilidad \parencite{KremerWillisYou2022} \vertema{3.A.45}. El PIB real por habitante de las economías emergentes y en desarrollo se ha multiplicado por 2,6 desde 1995, frente a 1,5 en las avanzadas, pero sobre todo por Asia \parencite{FMI2026WEO}.
\item \emph{La renta media}: muchos países que salieron de la pobreza se estancan al llegar a la renta media; es la \concepto{trampa de la renta media} \parencite{GillKharas2007East} (apartado~\ref{sec:formal}).
\item \emph{El desarrollo humano}: todas las regiones han mejorado su IDH desde 1990, y el África subsahariana sigue la última (gráfico~\ref{fig:idh-regiones}).
\end{itemize}

\begin{anotaciones}
\begin{figuracaja}
\caption{Índice de desarrollo humano por grupos de desarrollo humano y por regiones, 1990 y 2023}\label{fig:idh-regiones}
\grafico[0.8\textwidth]{idh-regiones}
\fuente{elaboración propia a partir de \textcite{PNUD2025Series} (Informe sobre Desarrollo Humano 2025). El índice y la lista de países, en \href{https://es.wikipedia.org/wiki/\%C3\%8Dndice_de_desarrollo_humano}{«Índice de desarrollo humano»} y \href{https://es.wikipedia.org/wiki/Anexo:Pa\%C3\%ADses_por_\%C3\%ADndice_de_desarrollo_humano}{«Anexo:Países por índice de desarrollo humano»} (Wikipedia).}
\end{figuracaja}
China, con un IDH de 0,797 (puesto 78), está ya en el grupo de desarrollo humano \emph{alto}, casi en el muy alto \parencite{PNUD2025IDH}.
\end{anotaciones}

\paragraph{Pobreza: una caída histórica, cada vez más concentrada en África}

\begin{itemize}
\item \emph{La evolución} (gráfico~\ref{fig:pobreza-regiones}): la pobreza extrema ha caído del 43,4~\% de la población mundial en 1990 al 10,1~\% en 2024, aunque desde 2014 lo hace muy despacio \parencite{BM2026PIP}.
  \begin{itemize}
  \item \emph{Asia}: la caída se concentra en Asia oriental (del 67~\% al 2~\%) y meridional (del 50~\% al 3~\%). China explica casi 2 tercios de la reducción mundial: sus pobres extremos pasaron de unos 940 millones en 1990, el 83~\% de su población, a prácticamente ninguno desde 2019 \parencite{BM2026PIP}.
  \item \emph{África}: en el África subsahariana, la tasa solo ha bajado del 62~\% al 45~\%, y como su población crece deprisa, el número de pobres ha aumentado de 321 a 581 millones: hoy son el 70~\% de los pobres extremos del mundo \parencite{BM2026PIP}.
  \end{itemize}
\end{itemize}

\begin{figure}[H]
\centering
\caption{Población en pobreza extrema por regiones, 1990-2026 (millones de personas que viven con menos de 3,00 dólares diarios de PPA de 2021)}\label{fig:pobreza-regiones}
\grafico[0.8\textwidth]{pobreza-regiones}
\fuente{elaboración propia a partir de \textcite{BM2026PIP} (Poverty and Inequality Platform, actualización del 22 de septiembre de 2026). 2025 y 2026, estimaciones adelantadas del Banco Mundial.}
\end{figure}

\paragraph{Desigualdad: ¿una fase del desarrollo?}

\begin{itemize}

  \item \emph{La curva de \autor{Kuznets}}: \textcite{Kuznets1955} propuso que, al desarrollarse un país, la desigualdad primero sube y luego baja, una \concepto{U invertida} entre la renta por habitante y el índice de \autor{Gini} (IG) \vertema{3.A.24}. Al principio, la población pasa de un sector agrario pobre e igualitario a uno urbano más rico y más desigual; cuando la mayoría ya está en el moderno, la educación y la redistribución la reducen. Es la misma dualidad de \autor{Lewis} (apartado~\ref{sec:etapa}).
  \item \emph{La evidencia no la confirma como ley}. En el corte transversal actual de 150 países apenas se ve: un ajuste cuadrático explica el 10~\% de las diferencias de desigualdad entre países, y la U que dibuja la empujan América Latina y el África austral, con Gini altos para su renta (gráfico~\ref{fig:kuznets}). Dentro de cada país, la U aparece rara vez \parencite{DeiningerSquire1998}. Y en las economías avanzadas la desigualdad ha vuelto a subir desde los 80: \textcite{Piketty2014} explica su caída de 1914-1945 por las guerras, la inflación y los impuestos progresivos, no por una ley del desarrollo \vertema{3.A.45}.
  \item \emph{La curva del elefante}: entre 1988 y 2008 las rentas que más crecieron en el mundo fueron las del centro de la distribución mundial (las clases medias de China y de otros emergentes) y las del 1~\% más rico; las que menos, las de las clases medias y bajas de los países ricos \parencite{LaknerMilanovic2016}. Como los emergentes crecieron más, la desigualdad \emph{entre} países ha caído desde los 90, aunque haya subido \emph{dentro} de muchos \parencite{Milanovic2016}.
\end{itemize}

\begin{figure}[H]
\centering
\caption{La curva de \autor{Kuznets} en los datos: índice de \autor{Gini} y PIB por habitante en 150 países (último año disponible entre 2015 y 2024)}\label{fig:kuznets}
\grafico[0.75\textwidth]{kuznets}
\fuente{elaboración propia a partir de \textcite{BM2026WDI}. PIB por habitante en dólares de PPA de 2021, en escala logarítmica; ajuste cuadrático en el logaritmo de la renta ($R^2=0{,}10$).}
\end{figure}

\begin{itemize}
\item \emph{Nexo}: unos países convergen y otros no, y la pobreza se concentra donde no ha habido crecimiento. ¿Por qué algunos países se quedan atrás? Es la pregunta de las teorías del desarrollo.
\end{itemize}

% --------------------------------------------------------------------------
\subsection{Las teorías pioneras y su replanteamiento}\label{sec:clasicas}

\begin{itemize}
\item \emph{Conclusión anticipada}: los primeros teóricos del desarrollo, de los 50 a los 70, vieron el atraso como una etapa (\autor{Lewis}, \autor{Rostow}) o como una situación estable, por causas internas (\autor{Myrdal}, \autor{Kaldor}, \autor{Nurkse}) o externas (la CEPAL); la escuela ortodoxa y el Consenso de Washington los replantearon desde el mercado.
\end{itemize}

\subsubsection{Emergencia e ideas comunes de la teoría del desarrollo}\label{sec:emergencia}

\paragraph{La insatisfacción con la teoría del crecimiento}


\begin{itemize}
\item \emph{Lo que decía la teoría del crecimiento}: las primeras teorías del desarrollo nacen a la vez que el modelo de \textcite{Solow1956} \vertema{3.A.43}, que llega a 2 resultados:
  \begin{itemize}
  \item los rendimientos decrecientes del capital llevan a un estado estacionario en el que solo el progreso técnico, exógeno, permite crecer de forma sostenida;
  \item los países con poco capital, en los que el capital rinde más, crecen más deprisa: hay \concepto{convergencia condicional}, la renta por habitante de cada país converge a su propio estado estacionario, y los pobres alcanzan a los ricos si tienen el mismo ahorro, la misma población y la misma tecnología \vertema{3.A.45}.
  \end{itemize}
\item \emph{Las 2 razones de la insatisfacción}:
  \begin{itemize}
  \item \emph{el modelo explicaba poco}: el progreso técnico, un residuo, daba cuenta del 87,5~\% del aumento de la producción por hora trabajada en Estados Unidos entre 1909 y 1949 \parencite{Solow1957}, la «medida de nuestra ignorancia» \parencite{Abramovitz1956};
  \item \emph{el mundo como nivelador}: la teoría neoclásica veía la economía mundial como un sistema que tiende a igualar las rentas, y los países pobres no parecían acercarse a los ricos.
  \end{itemize}
\end{itemize}

\paragraph{Las ideas comunes}\label{sec:ideas-comunes}


\begin{itemize}
\item \emph{Quiénes}: \autor{Lewis}, \autor{Rosenstein-Rodan}, \autor{Nurkse}, \autor{Myrdal}, \autor{Kaldor}, \autor{Hirschman}, \autor{Prebisch} y \autor{Singer}. Se les suele llamar heterodoxos, aunque su inspiración es más bien keynesiana: la economía puede estar en un equilibrio malo, la demanda importa y hace falta un impulso para salir de él.
\item \emph{Sus 2 ideas}:
  \begin{itemize}
  \item \emph{La convergencia no está garantizada}: las variables de las que depende el crecimiento en el modelo neoclásico (ahorro, capital humano, tamaño del mercado) dependen a su vez del desarrollo, y pueden formar círculos viciosos y equilibrios múltiples\footnote{La separación entre la teoría del desarrollo y la del crecimiento no fue tan tajante: muchas de estas ideas (economías de escala, externalidades, equilibrios múltiples) las incorporó después, ya formalizadas, la teoría del crecimiento endógeno (apartado~\ref{sec:formal}), con implicaciones de política no muy distintas.}.
  \item \emph{El cambio estructural es causa, y no solo consecuencia, del desarrollo}: los modelos neoclásicos tienen un solo sector, pero los países pobres son \emph{duales}, con un sector moderno y otro tradicional. Si el crecimiento agregado es la media de los de cada sector ponderada por su peso, trasladar trabajo de la agricultura poco productiva a la industria hace crecer la economía, y la industria es el motor del crecimiento.
  \end{itemize}
\item \emph{La excepción}: no todos pensaban así. \textcite{Schultz1964} sostenía que los campesinos tradicionales son pobres pero eficientes: asignan bien lo poco que tienen, y lo que les falta para crecer son factores nuevos (semillas mejoradas, fertilizantes, educación) que hagan rentable invertir. Es un germen de la contrarrevolución neoclásica (apartado~\ref{sec:ortodoxa}).
\item \emph{Nexo}: estas 2 ideas ordenan el resto del bloque. La dualidad lleva a \autor{Lewis} y a \autor{Rostow} (el subdesarrollo como etapa); la no convergencia, a \autor{Myrdal}, \autor{Kaldor}, \autor{Nurkse} y la CEPAL (el subdesarrollo como situación).
\end{itemize}

\subsubsection{El subdesarrollo entendido como una etapa}\label{sec:etapa}

\paragraph{Modelo dual de \autor{Lewis} (1954)}

\subparagraph{a) Idea}

\begin{itemize}
\item ¿Cómo crece una economía con un sector tradicional que tiene trabajo de sobra? \textcite{Lewis1954}\footnote{\autor{Arthur Lewis} recibió el Premio Nobel de Economía en 1979, con \autor{Theodore Schultz}, «por su investigación pionera sobre el desarrollo económico, con especial atención a los problemas de los países en desarrollo» \parencite[traducción propia]{NobelPrize1979}.} responde que trasladando ese trabajo, sin coste para el campo, a un sector moderno que reinvierte sus beneficios.
\item Para \autor{Lewis}, el problema central del desarrollo es entender cómo una economía que ahorra el 4 o el 5~\% de su renta pasa a ahorrar el 12 o el 15~\%: el ahorro sube porque sube la parte de la renta que va a quienes ahorran, los capitalistas.
\end{itemize}

\subparagraph{b) Modelo}

\begin{itemize}
\item \emph{i) Supuestos}:
  \begin{itemize}
  \item 2 sectores: el \emph{tradicional} (agrario, de subsistencia), donde la productividad marginal del trabajo es casi nula y se cobra un salario institucional, $\overline w_A$, el producto medio que reparte la familia; y el \emph{moderno} (industrial, capitalista), que maximiza beneficios;
  \item el sector moderno contrata a un salario $w_M=\overline w_A(1+\text{prima})$, con una prima que compensa el coste de la vida en la ciudad y del traslado; a ese salario la oferta de trabajo es \emph{perfectamente elástica}, la «oferta ilimitada de trabajo» del título del artículo;
  \item los capitalistas reinvierten sus beneficios; los trabajadores no ahorran.
  \end{itemize}
\item \emph{ii) Desarrollo}, paso a paso:
  \begin{itemize}
  \item \emph{la oferta de trabajo} (gráfico~\ref{fig:lewis-a}): el sector moderno puede contratar a cuantos quiera al salario $w_M$, por encima del salario institucional $\overline w_A$: la oferta es horizontal mientras haya excedente de trabajo en el campo, y solo se vuelve creciente al agotarse, en el \concepto{punto de \autor{Lewis}};
  \item \emph{el empleo y el beneficio} (gráfico~\ref{fig:lewis-b}): con el capital $K_1$, el empleo del sector moderno, $L_M$, se fija donde el valor del producto marginal del trabajo iguala al salario, $p\,\mathrm{PMg}_L(K_1)=w_M$: la demanda de trabajo es $D(K_1)$, el empleo es $L_1$ y el área entre la demanda y el salario es el beneficio de los capitalistas;
  \item \emph{la reinversión} (gráfico~\ref{fig:lewis}): el beneficio se reinvierte; con $K_2>K_1$ el producto marginal del trabajo sube, la demanda se desplaza a la derecha, $D(K_2)$, y el sector moderno contrata $L_2$ \emph{al mismo salario}, sin que caiga la producción agraria; y así sucesivamente con $K_3$, hasta $L_3$;
  \item \emph{el final del excedente}: en el punto de \autor{Lewis}, cada trabajador que se va reduce la producción del campo, los salarios suben en los 2 sectores y la economía pasa a funcionar como la neoclásica.
  \end{itemize}

\item \emph{iii) Implicaciones}:
  \begin{itemize}
  \item el motor del desarrollo es la \emph{acumulación del beneficio} del sector moderno, por el lado de la oferta: la parte de los beneficios en la renta y la tasa de ahorro suben a medida que la economía se industrializa;
  \item la desigualdad crece al principio, lo que enlaza con la curva de \autor{Kuznets} (apartado~\ref{sec:caracteristicas});
  \item la política debe favorecer la inversión del sector moderno y no subir los salarios antes de tiempo.
  \end{itemize}
\end{itemize}

\begin{figure}[H]
\centering
\caption{El modelo dual de \autor{Lewis} (a): la oferta ilimitada de trabajo}\label{fig:lewis-a}
\grafico[0.6\textwidth]{lewis-paso1}
\fuente{elaboración propia a partir de \textcite{Lewis1954} y \textcite[cap.~10]{Ray1998}. Primer paso del gráfico~\ref{fig:lewis}.}
\end{figure}

\begin{figure}[H]
\centering
\caption{El modelo dual de \autor{Lewis} (b): empleo y beneficio con el capital $K_1$}\label{fig:lewis-b}
\grafico[0.6\textwidth]{lewis-paso2}
\fuente{elaboración propia a partir de \textcite{Lewis1954} y \textcite[cap.~10]{Ray1998}. Segundo paso del gráfico~\ref{fig:lewis}.}
\end{figure}

\begin{figure}[H]
\centering
\caption{El modelo dual de \autor{Lewis} (c): la reinversión del beneficio desplaza la demanda de trabajo}\label{fig:lewis}
\grafico[0.6\textwidth]{lewis}
\fuente{elaboración propia a partir de \textcite{Lewis1954} y \textcite[cap.~10]{Ray1998}.}
\end{figure}


\subparagraph{c) Extensiones}

\begin{itemize}
\item \emph{\textcite{RanisFei1961}}: el trabajo del campo deja de sobrar antes de lo que parece. Cuando la productividad marginal agraria se hace positiva, cada trabajador que se va reduce la producción de alimentos, su precio relativo sube y con él el salario industrial, antes del punto de \autor{Lewis}. Industrializarse exige que crezca a la vez la productividad agraria: un crecimiento equilibrado entre los 2 sectores.
\item \emph{\textcite{HarrisTodaro1970}}: el paro está en las ciudades, no en el campo.\footnote{Un resumen del modelo, en \href{https://es.wikipedia.org/wiki/Modelo_de_Harris-Todaro}{«Modelo de Harris-Todaro»} (Wikipedia).} Si el salario urbano, $\overline w_M$, está fijado por encima del de equilibrio (por ley o por los sindicatos), los campesinos emigran hasta que el salario urbano \emph{esperado} iguala al del campo, $w_A$:
  \[
  w_A=\frac{L_M}{L_U}\,\overline w_M ,
  \]
  donde $L_U$ es la población activa urbana (los ocupados del sector moderno, $L_M$, más los parados) y $L_M/L_U$ la probabilidad de encontrar empleo en la ciudad\footnote{$L_U$ es una letra local de este párrafo. El modelo explica una paradoja de los países en desarrollo: se emigra a las ciudades aunque el paro urbano sea alto, porque lo que cuenta es el salario esperado y no el salario seguro.}. En equilibrio hay paro urbano, y crear empleo en la ciudad puede aumentarlo, porque atrae a más emigrantes del campo.
\end{itemize}

\subparagraph{d) Evidencia empírica}

\begin{itemize}
\item \emph{Las industrializaciones con salarios contenidos}: el modelo describe bien las de Japón, Corea del Sur o China, con décadas de salarios industriales casi estancados mientras el campo se vaciaba. En China, el debate sobre si se alcanzó el punto de \autor{Lewis} se abrió hacia 2010, cuando los salarios de los trabajadores emigrados empezaron a subir deprisa \parencite{ZhangYangWang2011}.
\item \emph{España}: el éxodo rural de los 60 dio a la industria una oferta de trabajo muy elástica que contuvo los salarios \vertema{4.A.27}.
\end{itemize}

\subparagraph{e) Valoración}

\begin{itemize}
\item \emph{A favor}: es el primer modelo que hace del cambio estructural el motor del crecimiento, y sigue siendo la forma de pensar la industrialización.
\item \emph{En contra}:
  \begin{itemize}
  \item los capitalistas pueden no reinvertir todo el beneficio, o sacarlo del país;
  \item la reinversión puede ahorrar trabajo (máquinas en lugar de empleo) y no desplazar la demanda de trabajo;
  \item supone una movilidad perfecta del trabajo entre sectores;
  \item el excedente de trabajo es menor de lo que suponía (\autor{Ranis} y \autor{Fei}), y el paro está en la ciudad (\autor{Harris} y \autor{Todaro}).
  \end{itemize}
\end{itemize}

\paragraph{Modelo de etapas de \autor{Rostow} (1960)}

\subparagraph{a) Idea}

\begin{itemize}
\item Si \autor{Lewis} explica el mecanismo, \textcite{Rostow1960} ordena la historia: todos los países pasan por las mismas etapas, y el subdesarrollo es solo la primera.
\end{itemize}

\subparagraph{b) Modelo}

\begin{itemize}
\item \emph{Las 5 etapas} (gráfico~\ref{fig:rostow}):
  \begin{enumerate}
  \item la \emph{sociedad tradicional}, agraria y con un techo tecnológico;
  \item las \emph{condiciones previas para el despegue}: inversión en transporte y en agricultura, un Estado nacional y una élite empresarial;
  \item el \emph{despegue}, la etapa clave: la inversión pasa de menos del 5~\% a más del 10~\% de la renta y unos sectores líderes arrastran al resto;
  \item el \emph{camino hacia la madurez}: la tecnología moderna se extiende a toda la economía;
  \item la \emph{era del alto consumo de masas}, con la terciarización.
  \end{enumerate}
\item \emph{La implicación: la brecha de financiación}. Para despegar hay que reunir el ahorro que pide la inversión, con ayuda exterior si hace falta. Los organismos internacionales lo calcularon con la versión de \autor{Domar} del modelo de \autor{Harrod} y \autor{Domar} \vertema{3.A.43}: si para crecer a una tasa objetivo $\bar\gamma_Y$ hace falta invertir
  \[
  \frac{I}{Y}=v\,\bar\gamma_Y ,
  \]
  con $v$ la relación capital-producto, la diferencia con la tasa de ahorro interna, $v\,\bar\gamma_Y-s$, era la ayuda «necesaria» \vertema{3.B.29}.
\end{itemize}

\begin{figure}[H]
\centering
\caption{Las etapas del crecimiento de \autor{Rostow}}\label{fig:rostow}
\grafico[0.6\textwidth]{rostow}
\fuente{elaboración propia a partir de \textcite{Rostow1960}.}
\end{figure}

\subparagraph{c) Valoración}

\begin{itemize}
\item \emph{Una descripción, más que una teoría}, inspirada en la historia británica. \textcite{Gerschenkron1962} mostró que los países atrasados no repiten las etapas: se las saltan, con más Estado o con más bancos para suplir lo que les falta.
\item \emph{La brecha de financiación no funcionó}: la ayuda no se tradujo en inversión ni la inversión en crecimiento, y el propio \autor{Domar} había renegado de ese uso de su modelo \parencites{Easterly1999}[cap.~2]{Easterly2003}.
\item \emph{Nada garantiza el despegue}: la idea de que el subdesarrollo es una etapa que acabará pasando llevó a otros autores a preguntarse lo contrario, por qué el atraso puede ser una situación estable.
\end{itemize}

\subsubsection{El subdesarrollo entendido como una situación: factores internos}\label{sec:internos}

\begin{itemize}
\item \emph{Nexo}: frente a la etapa, la situación. Estos autores buscan las fuerzas que perpetúan el atraso dentro de cada economía: la causación acumulativa, el crecimiento que se realimenta y las trampas de pobreza.
\end{itemize}

\paragraph{Efectos acumulativos: \autor{Myrdal} (1957)}

\subparagraph{a) Idea}

\begin{itemize}
\item El mercado no nivela, amplifica. Para \autor{Gunnar Myrdal}\footnote{\autor{Gunnar Myrdal} recibió el Premio Nobel de Economía en 1974, con \autor{Friedrich Hayek}, «por su trabajo pionero en la teoría del dinero y de las fluctuaciones económicas y por su penetrante análisis de la interdependencia de los fenómenos económicos, sociales e institucionales» \parencite[traducción propia]{NobelPrize1974}.}, la economía no tiende a un estado estacionario estable, como en el modelo de \autor{Solow}, sino a círculos virtuosos y viciosos de \concepto{causación circular acumulativa} \parencite{Myrdal1957}\footnote{El precedente es \textcite{Young1928Increasing}: las economías de escala dinámicas, que crecen con el tamaño del mercado, como motor del crecimiento.}.
\end{itemize}

\subparagraph{b) Modelo}

\begin{itemize}
\item \emph{El mecanismo}: la inversión va donde ya hay actividad, por las economías de aglomeración. La región rica atrae el capital y el trabajo cualificado de la pobre (\emph{efectos de retroceso}), y solo una parte de su prosperidad llega a la pobre (\emph{efectos de difusión}). Si dominan los primeros, las regiones ricas se hacen más ricas y las pobres, más pobres.
\item \emph{Dentro y fuera}: el mismo argumento vale entre las regiones de un país y entre países; por eso \autor{Myrdal} está entre los factores internos, donde actúa el mecanismo, y hace de puente con los externos (apartado~\ref{sec:externos}).
\item \emph{La implicación}: el Estado debe estimular las regiones a las que no va el mercado, para poner la causación en la senda virtuosa: es la lógica de la política regional, que formalizará la nueva geografía económica \vertema{3.B.6}.
\end{itemize}

\subparagraph{c) Valoración}

\begin{itemize}
\item Es una intuición poderosa, pero sin un modelo formal no dice cuándo domina el retroceso y cuándo la difusión. Esa pregunta es la que se hace \autor{Kaldor}.
\end{itemize}

\paragraph{Crecimiento acumulativo y equilibrios múltiples: \autor{Kaldor} (1966 y 1970)}

\subparagraph{a) Idea}

\begin{itemize}
\item \textcite{Kaldor1966,Kaldor1970} da forma a la causación acumulativa con la \concepto{ley de \autor{Kaldor}-\autor{Verdoorn}}: la productividad de la industria crece más cuanto más crece su producción (economías de escala dinámicas, aprendizaje por la práctica), y la producción crece más cuanto más productiva y competitiva es la economía.
\end{itemize}

\subparagraph{b) Modelo}

\begin{itemize}
\item \emph{2 ecuaciones}, que formalizan \textcite{DixonThirlwall1975}: el crecimiento de la productividad depende del de la producción, y el de la producción, a través de las exportaciones, del de la productividad.
\item \emph{El resultado}: todo depende de la fuerza de esa realimentación. Si es débil, la economía converge a una tasa de crecimiento que fijan las condiciones iniciales y la demanda externa; si es fuerte, se aleja de ella en una espiral virtuosa o viciosa, y pequeñas ventajas iniciales se perpetúan (anexo~\ref{anexo:clasicas}).
\item \emph{La implicación}: un impulso inicial (política industrial, promoción de las exportaciones) puede poner a un país en la buena senda.
\end{itemize}

\subparagraph{c) Valoración}

\begin{itemize}
\item \autor{Kaldor} contrastó la ley de \autor{Verdoorn} con los datos de la industria de los países de la OCDE, pero el modelo es de demanda y deja fuera la acumulación de factores. Su gran idea, que el estado estacionario único y estable puede no existir, es la de los equilibrios múltiples y las trampas de pobreza.
\end{itemize}

\paragraph{Trampas de pobreza: \autor{Nurkse} (1953)}

\subparagraph{a) Idea}

\begin{itemize}
\item Un país es pobre porque es pobre: para \textcite{Nurkse1953}, la pobreza se perpetúa por 2 \concepto{círculos viciosos}:
  \begin{itemize}
  \item por la \emph{oferta} de capital: renta baja $\to$ poco ahorro $\to$ poco capital $\to$ baja productividad $\to$ renta baja;
  \item por la \emph{demanda}: renta baja $\to$ mercado pequeño $\to$ pocos incentivos a invertir $\to$ poco capital $\to$ renta baja.
  \end{itemize}
\item El atraso es entonces un equilibrio estable de renta baja, una \concepto{trampa de pobreza}: pequeños aumentos de la renta se pierden y la economía vuelve a él.
\end{itemize}

\subparagraph{b) Modelo}

\begin{itemize}
\item \emph{i) Supuestos}: los del modelo de \autor{Solow}, pero con una función de producción con forma de S: rendimientos decrecientes con poco capital (una economía agraria), crecientes en un tramo intermedio (infraestructuras, externalidades o una tecnología moderna que solo compensa a partir de un umbral) y otra vez decrecientes con mucho \parencite{AzariadisDrazen1990}. Detrás puede haber un ahorro que depende de la renta, mercados financieros que no prestan a los pobres, impaciencia por la inseguridad o una población que crece más cuando la renta es baja\footnote{Una versión es la \emph{trampa demográfica} de \textcite{Nelson1956}: con rentas bajas, cada aumento de la renta eleva el crecimiento de la población más que el de la producción, y la renta por habitante vuelve a su nivel de subsistencia (el punto $A$); solo un salto que la lleve más allá del umbral $U$, a partir del cual la producción crece más que la población, saca al país de la trampa. Para conseguirlo, más ahorro, menos crecimiento de la población o un gran empujón.\graficonota{trampa-demografica}}.
\item \emph{ii) Desarrollo}, paso a paso, en la representación en tasas del modelo de \autor{Solow} \vertema{3.A.43}:
  \begin{itemize}
  \item \emph{las 2 curvas} (gráfico~\ref{fig:trampa-pobreza-a}): con la función en S, la curva de ahorro, $s\,f(\tilde k)/\tilde k$, ya no es monótona: decrece, crece y vuelve a decrecer; la de depreciación efectiva, $n+g+\delta$, sigue siendo horizontal;
  \item \emph{los 3 estados estacionarios} (gráfico~\ref{fig:trampa-pobreza-b}): las 2 curvas se cortan 3 veces. Donde la curva de ahorro está por encima de la horizontal, $\tilde k$ crece; donde está por debajo, decrece. Por eso hay 2 estados estacionarios estables, uno pobre ($\tilde k_B$, la trampa) y otro rico ($\tilde k_A$), separados por uno inestable ($\tilde k_U$), el umbral;
  \item \emph{la salida} (gráfico~\ref{fig:trampa-pobreza}): una economía en $\tilde k_B$ vuelve a $\tilde k_B$ tras un aumento pequeño del capital; solo uno que la lleve más allá de $\tilde k_U$ la pone en camino de $\tilde k_A$.
  \end{itemize}
\item \emph{iii) Implicaciones}: la renta a largo plazo depende de dónde se parte. Una ayuda pequeña no saca de la trampa; un \concepto{gran empujón} (\emph{big push}) que lleve el capital más allá de $\tilde k_U$, sí, y después puede retirarse.
\end{itemize}

\begin{figure}[H]
\centering
\caption{Una trampa de pobreza en el modelo de \autor{Solow} (a): la curva de ahorro con la función en S}\label{fig:trampa-pobreza-a}
\grafico[0.55\textwidth]{trampa-pobreza-paso1}
\fuente{elaboración propia a partir de \textcite[cap.~1]{SalaiMartin2000} y \textcite[cap.~1]{BarroSalaiMartin2004}. Primer paso del gráfico~\ref{fig:trampa-pobreza}.}
\end{figure}

\begin{figure}[H]
\centering
\caption{Una trampa de pobreza en el modelo de \autor{Solow} (b): 2 estados estacionarios estables y un umbral}\label{fig:trampa-pobreza-b}
\grafico[0.55\textwidth]{trampa-pobreza-paso2}
\fuente{elaboración propia a partir de \textcite[cap.~1]{SalaiMartin2000} y \textcite[cap.~1]{BarroSalaiMartin2004}. Segundo paso del gráfico~\ref{fig:trampa-pobreza}.}
\end{figure}

\begin{figure}[H]
\centering
\caption{Una trampa de pobreza en el modelo de \autor{Solow} (c): el gran empujón}\label{fig:trampa-pobreza}
\grafico[0.55\textwidth]{trampa-pobreza}
\fuente{elaboración propia a partir de \textcite[cap.~1]{SalaiMartin2000} y \textcite[cap.~1]{BarroSalaiMartin2004}.}
\end{figure}


\subparagraph{c) Extensiones: el gran empujón}

\begin{itemize}
\item \emph{La idea de \textcite{RosensteinRodan1943Problems}}: una fábrica de zapatos que emplea a miles de campesinos no es rentable sola, porque sus trabajadores no gastarán todo su salario en zapatos; sí lo es si a la vez se abren muchas fábricas de bienes distintos, porque los salarios de cada una compran lo que producen las demás. Es un \concepto{fallo de coordinación}: invertir en cada sector solo es rentable si invierten los demás.
  \begin{itemize}
  \item «Una golondrina no hace verano»: una fábrica sola no saca a un país de la pobreza; hace falta un programa de inversión simultáneo en muchos sectores, que coordina el Estado (el \emph{crecimiento equilibrado} que pedía también \autor{Nurkse}).
  \end{itemize}
\item \emph{La formalización de \textcite{MurphyShleiferVishny1989Industrialization}}: con 2 técnicas en cada sector, la tradicional y una moderna con un coste fijo y una prima salarial, una empresa moderna sola vende poco y no cubre el coste fijo; si se modernizan todas, los salarios y beneficios que pagan aumentan la demanda de todos los bienes y cada una vende lo bastante. Hay 2 equilibrios, todos tradicionales y todos modernos; el segundo es mejor para todos, pero el mercado puede quedarse en el primero (el álgebra y su gráfico, en el anexo~\ref{anexo:clasicas}).
\item \emph{La alternativa de \textcite{Hirschman1958}}: un país pobre no puede financiar un programa en todos los sectores a la vez. Mejor un crecimiento \emph{desequilibrado}: concentrar la inversión en los sectores con más \concepto{eslabonamientos}, hacia atrás (los que compran mucho a otros, como el automóvil) y hacia delante (los que venden a muchos, como el acero), para que sus cuellos de botella tiren de los demás. Se miden con las tablas \emph{input-output} \vertema{3.A.31}.
\end{itemize}

\subparagraph{d) Evidencia empírica}

\begin{itemize}
\item \emph{A escala de países, pocas}: casi todos los países pobres han crecido a largo plazo, y la renta de los más pobres no se queda quieta, como haría en una trampa \parencite{KraayMcKenzie2014Do}. Los mecanismos que la generarían (nutrición, ahorro, demografía) tampoco se observan con la fuerza necesaria \parencite[cap.~5]{Campante2021}.
\item \emph{A escala de hogares, más plausibles}: con activos o capacidades por debajo de un umbral, un hogar puede quedarse atrapado aunque el país crezca (apartado~\ref{sec:cultural}).
\end{itemize}

\subparagraph{e) Valoración}

\begin{itemize}
\item La trampa no la predijo \autor{Solow}: es de la literatura posterior, que cambió la forma de la función de producción. Su lección es que la política puede importar mucho en el punto de partida.
\item Pero justificó grandes programas de inversión pública y de ayuda que a menudo fracasaron: si la trampa es rara entre países, el empujón no basta, y hay que buscar qué falla en cada caso. La pregunta vuelve, formalizada, en las teorías recientes (apartado~\ref{sec:formal}).
\end{itemize}

\subsubsection{El subdesarrollo entendido como una situación: factores externos}\label{sec:externos}

\begin{itemize}
\item \emph{Nexo}: las teorías anteriores buscan las causas del atraso dentro de cada país; las siguientes, en sus relaciones con los países ricos.
\end{itemize}

\paragraph{La escuela estructuralista de la CEPAL: \autor{Prebisch} y \autor{Singer}}

\subparagraph{a) Idea}

\begin{itemize}
\item \emph{La correa de transmisión rota}: en el siglo XIX, el comercio transmitió el crecimiento del centro a la periferia (demanda de materias primas, capital y emigrantes); en el XX dejó de hacerlo. Para la Comisión Económica para América Latina (CEPAL), la economía mundial se divide en un \emph{centro} industrial y una \emph{periferia} que exporta productos primarios, y el intercambio entre ambos perjudica a la periferia \parencite{Prebisch1949,Singer1950}.\footnote{Un mapa del centro y la periferia del sistema mundial, en \href{https://elordenmundial.com/mapas-y-graficos/la-teoria-del-centro-periferia/}{«La teoría del centro-periferia»} (El Orden Mundial). El mismo enfoque aplicado a la Unión Europea, en \textcite{DelRio2018CentroPeriferia}.}
\end{itemize}

\subparagraph{b) Modelo}

\begin{itemize}
\item \emph{La tesis de \autor{Prebisch} y \autor{Singer}}: la relación real de intercambio de la periferia, $\mathrm{RRI}=p_X/p_M$, tiende a empeorar a largo plazo:
  \begin{itemize}
  \item porque la demanda de productos primarios crece menos que la renta (ley de \autor{Engel}) y el centro protege su agricultura;
  \item porque en el centro los sindicatos y el poder de mercado retienen las ganancias de productividad en forma de salarios y márgenes, mientras que en la periferia la competencia las traslada a los precios.
  \end{itemize}
  Cada unidad importada le cuesta a la periferia cada vez más exportaciones: el fruto de su productividad se lo queda, a través de los precios, el centro.
\item \emph{La implicación}: la \concepto{industrialización por sustitución de importaciones} (ISI), proteger la industria nacional hasta que aprenda a producir con costes competitivos, con el argumento de la industria naciente \vertema{3.B.8}\footnote{La ISI tenía también un pilar monetario: tras la Segunda Guerra Mundial los países pobres no tenían divisas para importar, y producir en casa lo que se importaba era una forma de racionarlas \parencite{Irwin2021}. \textbf{Para ver.} Hoy defiende la protección de la industria naciente \autor{Ha-Joon Chang}, en \href{https://www.youtube.com/watch?v=mi49GjJOPm0}{«Why Are Some Countries Rich and Others Poor? | Economics for People with Ha-Joon Chang»} (New Economic Thinking, YouTube).}. Se aplicó en América Latina, la India, buena parte de África o España, de los 50 a los 70.
\item \emph{Una versión moderna}: la restricción externa al crecimiento, la \emph{ley de \autor{Thirlwall}}\footnote{A largo plazo, un país no puede crecer más deprisa de lo que permite su balanza de pagos: $\gamma_Y=(\varepsilon_{X,Y^*}/\varepsilon_{M,Y})\,\gamma_{Y^*}$, donde $\varepsilon_{X,Y^*}$ es la elasticidad de sus exportaciones a la renta mundial, $\varepsilon_{M,Y}$ la de sus importaciones a su renta y $\gamma_{Y^*}$ el crecimiento del resto del mundo \parencite{Thirlwall1979}. Si la periferia exporta bienes con poca elasticidad renta e importa bienes con mucha, crece menos que el centro.}, y, en una corriente más radical, la \emph{teoría de la dependencia}, que atribuye el subdesarrollo a la propia integración en el capitalismo mundial \parencite{Frank1967,CardosoFaletto1969}.
\end{itemize}

\subparagraph{c) Valoración}

\begin{itemize}
\item \emph{La tendencia}: la caída de la relación real de intercambio de los productos primarios se observa, pero es débil y con rupturas.
\item \emph{La ISI}: creó industria, pero cara e ineficiente, con búsqueda de rentas en el reparto de las licencias \parencite{Krueger1974Political}, crisis de balanza de pagos y pocas exportaciones; la India era en 1980 más pobre en relación con Estados Unidos que en 1950 (el 5~\% de su renta por habitante, frente al 6,5~\%) \parencite{BoltvanZanden2025}. A favor, la historia: los países ricos se industrializaron con la protección que después negaron a los pobres \parencite{Chang2002}.
\item \emph{Hoy vuelve}, como política industrial de autonomía estratégica, de \emph{Made in China 2025} a los planes de la India \parencite{Irwin2021}. La evidencia y la comparación con la alternativa asiática se estudian aparte \vertema{3.B.9}.
\end{itemize}

\paragraph{La alternativa: la industrialización orientada a la exportación}

\begin{itemize}
\item \emph{La idea}: la que siguieron desde los 60 los \emph{4 tigres asiáticos} (Hong Kong, Singapur, Taiwán y Corea del Sur): industrializarse vendiendo al mundo, empezando por las manufacturas intensivas en trabajo en las que tenían ventaja comparativa y subiendo de gama a medida que acumulaban capital humano.
\item \emph{Los instrumentos}: no solo subvenciones a la exportación, sino un tipo de cambio realista, acceso libre a los insumos importados para exportar y apoyos (crédito, subvenciones) condicionados a vender en los mercados exteriores, que disciplinaban a las empresas.
\item \emph{La valoración}: crecimientos del 6 o el 7~\% anual durante décadas. Pero no fue solo el comercio: también contaron el ahorro alto, la educación, la estabilidad macroeconómica, la reforma agraria previa en Corea del Sur y Taiwán y la inversión extranjera (anexo~\ref{anexo:clasicas}) \vertema{3.B.9}.
\end{itemize}

% --------------------------------------------------------------------------
\subsubsection{Replanteamiento: la escuela ortodoxa y el Consenso de Washington}\label{sec:ortodoxa}

\begin{itemize}
\item \emph{Nexo}: mientras unos protegían y otros exportaban, la escuela ortodoxa recuperó el pensamiento neoclásico para aplicarlo al subdesarrollo. Su historia va de los 60 a los 90, en 4 pasos.
\end{itemize}

\paragraph{De la reconsideración conservadora a la contrarrevolución neoclásica}

\begin{itemize}
\item \emph{Los 60 y 70: la reconsideración conservadora}. \autor{Bauer} y \autor{Harry Johnson} critican los 3 pilares de las teorías anteriores: la protección (el libre comercio da ganancias), la industrialización a toda costa (la agricultura también puede ser motor) y el intervencionismo (los gobiernos también fallan: búsqueda de rentas, grupos de presión). Para \textcite{Bauer1972}, la ayuda al desarrollo refuerza al Estado frente al mercado y puede perjudicar.
\item \emph{Los 70: el giro hacia las necesidades básicas}. Países como Brasil crecían sin reducir la pobreza. La Conferencia Mundial del Empleo de la Organización Internacional del Trabajo (Ginebra, 1976) propuso orientar el desarrollo a las \concepto{necesidades básicas} \parencite{OIT1976NecesidadesBasicas}:
  \begin{itemize}
  \item lo imprescindible para cada familia: comer, vestirse y tener un techo;
  \item servicios esenciales: agua potable, saneamiento, salud, educación y transporte;
  \item un empleo adecuadamente remunerado;
  \item la participación en las decisiones que afectan a cada uno.
  \end{itemize}
\item \emph{Los 80: la contrarrevolución neoclásica}. 3 hechos la impulsan:
  \begin{itemize}
  \item la crisis de la deuda de 1982, tras una década de préstamos dedicados en buena parte al consumo;
  \item el agotamiento de la ISI, cuyos países no convergían;
  \item el éxito de los tigres asiáticos, leído como el triunfo del mercado y de la apertura.
  \end{itemize}
  \autor{Krueger}, \autor{Bhagwati} o \autor{Balassa} defienden que el desarrollo lo traen la apertura y los precios correctos, y los programas de ajuste del Fondo Monetario Internacional (FMI) y del Banco Mundial lo aplican en los países endeudados.
\end{itemize}

\paragraph{El Consenso de Washington}

\subparagraph{a) Idea}

\begin{itemize}
\item \textcite{Williamson1990} resumió en una conferencia de noviembre de 1989, en 10 puntos, lo que los organismos con sede en Washington (el FMI, el Banco Mundial, el Tesoro de Estados Unidos) pedían a América Latina: es el \concepto{Consenso de Washington}.
\end{itemize}

\subparagraph{b) Modelo}

\begin{itemize}
\item \emph{Sus 4 bloques} (los 10 puntos, en el anexo~\ref{anexo:clasicas}):
  \begin{itemize}
  \item \emph{estabilidad macroeconómica}: disciplina fiscal, gasto público reorientado a la sanidad, la educación y las infraestructuras, y una reforma tributaria con bases amplias y tipos moderados;
  \item \emph{precios de mercado}: tipos de interés libres y un tipo de cambio competitivo;
  \item \emph{apertura}: liberalización comercial y de la inversión directa extranjera;
  \item \emph{mercado y Estado}: privatización, desregulación y derechos de propiedad seguros.
  \end{itemize}
\end{itemize}

\subparagraph{c) Valoración}

\begin{itemize}
\item \emph{La crisis del Consenso}: en los 90, América Latina creció poco, las crisis financieras (México, 1994; Asia, 1997; Rusia, 1998; Argentina, 2001) mostraron el riesgo de liberalizar las entradas de capital sin regulación, y la transición de Rusia fracasó. Las críticas:
  \begin{itemize}
  \item \emph{«café para todos»}: la misma receta para países distintos, sin tener en cuenta sus instituciones ni el orden de las reformas (primero la cuenta corriente y la inversión directa, después la de cartera);
  \item \emph{una lectura parcial de Asia}: el informe del Banco Mundial sobre el «milagro» de Asia oriental \parencite{BM1993Milagro} y \textcite{Stiglitz1996} reconocen que Corea del Sur, Taiwán o Singapur no habían liberalizado sin más; habían intervenido (crédito dirigido, promoción selectiva de exportaciones), aunque para disciplinar al mercado y no para sustituirlo;
  \item \emph{el segundo óptimo}: en un mundo con fallos de mercado y de gobierno que no se pueden eliminar a la vez, copiar las «mejores prácticas» de los ricos puede empeorar las cosas \vertema{3.A.22} (apartado~\ref{sec:institucional});
  \item \emph{la equidad}: el ajuste olvidó los costes sociales a corto plazo.
  \end{itemize}
\item \emph{Las correcciones}: el Banco Mundial ya había propuesto en 1991 un enfoque \emph{favorable al mercado} (\emph{market-friendly}), en el que el Estado invierte en las personas y en infraestructuras y deja el resto al mercado \parencite{BM1991Informe}; en 1998 reconoció que hacían falta reformas institucionales \parencite{BurkiPerry1998}, y su presidente propuso un \emph{Consenso de Santiago}, con las reformas de segunda generación: educación, instituciones, sector financiero y pobreza \parencite{Wolfensohn1998}.
\item \emph{Lo que queda}: pese a su desprestigio, la dimensión macroeconómica del Consenso (disciplina fiscal, inflación baja, tipo de cambio realista) se ha consolidado como condición necesaria del desarrollo, aunque no suficiente.
\item \emph{Recapitulación y nexo}: los teóricos del desarrollo vieron que el atraso puede ser un equilibrio estable (círculos viciosos, fallos de coordinación, restricción externa), y sus recetas (el gran empujón, la ISI) y las de la contrarrevolución (el Consenso de Washington) fallaron por lo mismo: eran recetas generales para problemas particulares. Desde finales de los 80, la teoría formaliza esas intuiciones y se pregunta por las instituciones.
\end{itemize}

% --------------------------------------------------------------------------
\subsection{Teorías recientes: complementariedades e instituciones}\label{sec:recientes}

\begin{itemize}
\item \emph{Conclusión anticipada}: las teorías recientes convierten las intuiciones de los pioneros (rendimientos crecientes, externalidades, equilibrios múltiples) en modelos que se pueden contrastar, y buscan la causa última del atraso en las instituciones.
\end{itemize}

\subsubsection{La teoría formal: crecimiento endógeno, fallos de coordinación y complementariedades}\label{sec:formal}


\paragraph{El crecimiento endógeno}

\begin{itemize}
\item \emph{La idea}: si el capital, en sentido amplio (físico y humano), no tiene rendimientos decrecientes, como en el modelo $AK$, $Y=AK$, la economía crece a una tasa que depende del ahorro,
  \[
  \gamma_{Y/L}=s\,A-(n+\delta),
  \]
  y esa tasa no cae aunque el país se enriquezca \parencite{Rebelo1991Long} \vertema{3.A.44}.
\item \emph{Las 2 ideas que los teóricos del desarrollo habían adelantado}:
  \begin{itemize}
  \item \emph{la convergencia no está garantizada}: si la rentabilidad no cae con la renta, los pobres no tienen por qué crecer más que los ricos;
  \item \emph{la economía se aleja de la competencia perfecta}: con externalidades o con innovación que exige poder de mercado, el equilibrio no es eficiente y hay espacio para la política.
  \end{itemize}
\item \emph{Por qué el capital no va a los pobres}: con rendimientos decrecientes, el capital debería rendir muchísimo más en la India que en Estados Unidos y fluir de los ricos a los pobres; no lo hace, y es la \emph{paradoja de \autor{Lucas}} \parencite{Lucas1990Why}\footnote{\autor{Lucas} calcula que, con la misma tecnología y sin diferencias de capital humano, el producto marginal del capital en la India sería decenas de veces el de Estados Unidos. Las diferencias de capital humano, de productividad, de instituciones y de riesgo explican por qué no es así \vertema{3.B.29}.}.
\end{itemize}

\paragraph{Los fallos de coordinación}

\begin{itemize}
\item Volvamos a la trampa de pobreza (gráfico~\ref{fig:trampa-pobreza}): el gran empujón de \autor{Rosenstein-Rodan}, que \textcite{MurphyShleiferVishny1989Industrialization} formalizan con rendimientos crecientes, es un \concepto{fallo de coordinación}: una técnica moderna solo es rentable si se vende lo bastante, y eso solo ocurre si los demás sectores también se modernizan y pagan salarios que se gastan en sus productos. En un equilibrio nadie invierte porque nadie invierte; en el otro, todos (anexo~\ref{anexo:clasicas}).
\item Lo que en \autor{Nurkse} era una intuición es ahora un resultado: el mercado puede quedarse en el equilibrio malo aunque el bueno sea mejor para todos, y salir de él exige coordinar expectativas e inversiones.
\end{itemize}

\paragraph{Las complementariedades: la teoría del anillo O}

\begin{itemize}
\item \emph{La idea}: el transbordador \emph{Challenger} estalló en 1986 porque falló una junta, un anillo de goma barato. \textcite{Kremer1993ORing} lo convierte en una teoría del desarrollo: muchos procesos productivos son una cadena de tareas, y basta que una salga mal para que se pierda casi todo el valor.
\item \emph{El modelo}: un producto exige $m$ tareas, y cada una sale bien con una probabilidad $q_i$ que depende de la cualificación de quien la hace. Con $\overline V$ el valor del producto si todo sale bien ($m$, $q_i$ y $\overline V$ son letras locales), el valor esperado por trabajador es proporcional a
  \[
  \overline V\prod_{i=1}^{m}q_i .
  \]
  Como las cualificaciones se multiplican, pequeñas diferencias dan grandes diferencias: con 10 tareas, si cada una sale bien el 95~\% de las veces se obtiene el 60~\% del valor; si sale bien el 90~\%, solo el 35~\%.
\item \emph{Las implicaciones}:
  \begin{itemize}
  \item los trabajadores se emparejan por cualificación, porque los buenos rinden más con los buenos, y los salarios crecen más que proporcionalmente con ella: los cualificados de los países pobres ganan mucho más si emigran (la fuga de cerebros);
  \item los países con poca cualificación se especializan en productos sencillos, con pocas tareas y poco valor añadido;
  \item no formarse es racional si nadie más lo hace: otra trampa por complementariedades. Y unas infraestructuras malas (cortes de luz, transporte) son un anillo defectuoso para toda la economía.
  \end{itemize}
\end{itemize}

\paragraph{La distancia a la frontera y la trampa de la renta media}

\begin{itemize}
\item \emph{La idea}: lo que hace crecer a un país depende de lo lejos que esté de la frontera tecnológica \parencite{AcemogluAghionZilibotti2006}.
  \begin{itemize}
  \item Lejos de ella, crecer \emph{imitando} es eficaz, y convienen instituciones que favorecen la inversión a gran escala: grandes grupos empresariales, crédito bancario, relaciones a largo plazo.
  \item Cerca de ella hay que \emph{innovar}, y convienen otras: competencia, entrada de empresas, mercados de capitales, universidades.
  \item Un país que no cambia de instituciones a tiempo deja de converger \vertema{3.A.44}.
  \end{itemize}
\item \emph{La \concepto{trampa de la renta media}} \parencite{GillKharas2007East}: muchos países crecen deprisa hasta la renta media y allí se estancan, cuando se agotan la ventaja de los salarios bajos y la imitación y aún no saben innovar. Desde 1990 solo 34 economías de renta media han llegado a renta alta, más de un tercio gracias a la integración en la Unión Europea o al petróleo, y 108 siguen en la renta media, con unos 6.000 millones de personas \parencite{BM2024Informe}.
\item \emph{La receta del Banco Mundial}: invertir cuando se es pobre, \emph{incorporar} tecnología de fuera al llegar a la renta media-baja e \emph{innovar} en la media-alta, como hizo Corea del Sur \parencite{BM2024Informe}.
\item \emph{Recapitulación y nexo}: los modelos formales explican por qué puede no haber convergencia, pero no por qué unos países coordinan, acumulan e innovan y otros no. La respuesta más robusta mira a las reglas del juego.
\end{itemize}

\subsubsection{Enfoque institucional: la importancia de las instituciones}\label{sec:institucional}


\begin{itemize}
\item \emph{Conclusión anticipada}: las instituciones son la causa última más robusta del desarrollo, pero saber que importan no dice cuáles hay que tener ni cómo cambiarlas.
\item \emph{Qué son}: las \concepto{instituciones} son las reglas del juego de una sociedad \parencite[p.~3]{North1990}\footnote{\autor{Douglass North} recibió el Premio Nobel de Economía en 1993, con \autor{Robert Fogel}, «por haber renovado la investigación en historia económica aplicando la teoría económica y métodos cuantitativos para explicar el cambio económico e institucional» \parencite[traducción propia]{NobelPrize1993North}.}: la protección de la propiedad, el cumplimiento de los contratos, los límites al poder, la facilidad para crear una empresa. Importan porque fijan los incentivos a invertir, innovar y comerciar, y porque deciden quién se queda con los frutos.
\item \emph{El contexto}: tras el Consenso de Washington, el propio Banco Mundial admitió que las instituciones importan \parencite{BurkiPerry1998}, y \textcite{HallJones1999Why} atribuyeron las diferencias de productividad entre países a su infraestructura social.
\item \emph{Inclusivas y extractivas}: \textcite{AcemogluRobinson2012} distinguen las instituciones \concepto{inclusivas}, que protegen la propiedad de todos, abren los mercados y reparten el poder, de las \concepto{extractivas}, que concentran la renta y el poder en una élite.
\item \emph{La evidencia, en una frase}: \textcite{AcemogluJohnsonRobinson2001Colonial} usan la mortalidad de los colonos europeos como instrumento, correlacionado con las instituciones y con el \emph{nivel} de la renta actual, no con su crecimiento: el efecto estimado con variables instrumentales (0,94) es mayor que el de mínimos cuadrados (0,52) \vertema{3.A.45}.
\item \emph{Por qué persisten las malas}: porque benefician a quien tiene el poder, que perdería con el cambio; es la visión del conflicto social, la que defienden \textcite{AcemogluJohnsonRobinson2005} \vertema{3.A.45}.
\item \emph{Un ejemplo}: Botsuana, uno de los países más pobres del mundo al independizarse en 1966, conservó instituciones tradicionales que limitaban el poder de los jefes, protegió la propiedad y usó la renta de los diamantes sin caer en la enfermedad holandesa (apartado~\ref{sec:caracteristicas}): desde entonces ha sido de los países que más han crecido del mundo \parencite{AcemogluJohnsonRobinson2003Botswana}.
\item \emph{Qué son buenas instituciones: no hay consenso}. La misma función (proteger la propiedad, dar crédito) se cumple con formas muy distintas: China creció en los 80 con empresas colectivas de los municipios y un sistema de doble vía de precios que no habría recomendado ningún manual. \textcite{Rodrik2008Second} aconseja \concepto{instituciones de segundo óptimo}, adaptadas a los fallos de mercado y de gobierno que cada país no puede eliminar a corto plazo, en lugar de copiar las «mejores prácticas» de los ricos \vertema{3.A.22}.
  \begin{itemize}
  \item «El hábito no hace al monje»: lo que importa es la función de la institución, no su forma.
  \end{itemize}
\item \emph{El diagnóstico del crecimiento}: si no hay una receta para todos, ¿por dónde empezar? \textcite{HausmannRodrikVelasco2008} proponen buscar la restricción que más limita el crecimiento de cada país, la que tiene el precio sombra más alto, y concentrar ahí las reformas. Si la inversión privada es baja, es:
  \begin{itemize}
  \item o porque su \emph{rendimiento social} es bajo (geografía, capital humano, infraestructuras);
  \item o porque el inversor no puede \emph{apropiárselo} (impuestos, corrupción, inseguridad jurídica, inestabilidad macroeconómica, fallos de información y de coordinación);
  \item o porque \emph{financiarse} es caro (poco ahorro, mala intermediación).
  \end{itemize}
  En cada rama se busca la señal de la restricción que aprieta: si es la financiación, los tipos de interés reales serán altos y los proyectos rentables no encontrarán crédito.
\item \emph{El desarrollo como autodescubrimiento}: descubrir qué puede producir bien un país es caro, y quien lo descubre no se queda con el beneficio, porque los demás lo imitan; se descubre demasiado poco. La política industrial debe premiar el descubrimiento y retirar el apoyo si no hay resultados \parencite{HausmannRodrik2003,Rodrik2004Industrial}.
\item \emph{Recapitulación y nexo}: las instituciones explican mucho, pero cambian despacio. Mientras tanto, ¿qué se puede hacer, dadas las instituciones heredadas, para que a los pobres les vaya mejor? Es la pregunta de la aproximación microeconómica.
\end{itemize}

% ==========================================================================
\section{Aproximación microeconómica}\label{sec:microap}
\minutos{13,5}

\begin{notaopositor}
En el cante (\tiempo{13,5}): las contribuciones pioneras de \autor{Sen} y \autor{Deaton} (\tiempo{2}); el comportamiento y los mercados de los pobres (\tiempo{3}): aspiraciones, escasez y trampas de los hogares (\tiempo{1,5}), y crédito y seguros (\tiempo{1,5}); la aproximación experimental (\tiempo{6}): el método (\tiempo{2}) con la descomposición del sesgo de selección \eqref{eq:sesgo-seleccion}, la evidencia (\tiempo{2,5}) con 4 o 5 ejemplos y la valoración (\tiempo{1,5}); las implicaciones para las políticas (\tiempo{2,5}), con las 6 lecciones. La tabla de experimentos (anexo~\ref{anexo:experimentos}), la demostración del anexo~\ref{anexo:evaluacion} y las notas son de profundización.
\end{notaopositor}

\begin{itemize}
\item \emph{Conclusión anticipada}: a escala de las personas, la pobreza no solo es falta de renta: cambia cómo se decide y deja a los hogares fuera de mercados que funcionan mal. La aproximación experimental ha dado resultados sólidos y a veces sorprendentes sobre qué programas funcionan, aunque no sustituye a la teoría ni a la macroeconomía.
\item \emph{La idea}: en lugar de buscar la causa agregada del atraso, estudiar cómo viven y deciden los pobres y cómo funcionan los mercados concretos (crédito, seguros, educación, salud). Si \autor{Acemoglu}, \autor{Johnson} y \autor{Robinson} preguntan cómo la historia formó las instituciones, \autor{Duflo} pregunta qué se puede hacer, dadas esas instituciones, para estimular el desarrollo.
\end{itemize}

% --------------------------------------------------------------------------
\subsection{Las contribuciones pioneras: \autor{Sen} y \autor{Deaton}}\label{sec:pioneras}

\begin{itemize}
\item \emph{El precedente}: \textcite{Schultz1964} ya había mirado al campesino y no al agregado: los campesinos tradicionales son pobres pero eficientes, y asignan bien lo poco que tienen (apartado~\ref{sec:ideas-comunes}).
\item \emph{\autor{Amartya Sen}}\footnote{\autor{Amartya Sen} recibió el Premio Nobel de Economía en 1998 «por sus contribuciones a la economía del bienestar» \parencite[traducción propia]{NobelPrize1998}.}: cambia la pregunta de cuánto se tiene a qué se puede hacer con ello.
  \begin{itemize}
  \item \emph{Las capacidades}: el bienestar no está en los bienes, sino en lo que permiten ser y hacer (estar sano, saber leer, participar); la pobreza es privación de capacidades \parencite{Sen1985Commodities,Sen1999}. Es el fundamento del IDH y del índice de pobreza multidimensional (apartado~\ref{sec:medicion}).
  \item \emph{Las titularidades y las hambrunas}: la hambruna de Bengala de 1943 se produjo sin que cayera la producción de alimentos. Lo que se hundió fue lo que los jornaleros podían obtener con su trabajo, sus \concepto{titularidades} (\emph{entitlements}), porque los precios subieron y sus salarios no \parencite{Sen1981Poverty}. El hambre es un problema de acceso, no solo de disponibilidad, y las democracias con prensa libre, que obligan a los gobiernos a reaccionar, no han tenido grandes hambrunas \parencite{Sen1999}.
  \end{itemize}
\item \emph{\autor{Angus Deaton}}\footnote{\autor{Angus Deaton} recibió el Premio Nobel de Economía en 2015 «por su análisis del consumo, la pobreza y el bienestar» \parencite[traducción propia]{NobelPrize2015}.}: lleva la microeconometría al estudio de la pobreza.
  \begin{itemize}
  \item \emph{Las encuestas de hogares}: medir el consumo de cada hogar, y no solo la renta nacional, para saber quién es pobre, cómo reacciona a los precios y cómo le afecta una política (un impuesto sobre el arroz, una subvención) \parencite{Deaton1997Analysis}. Las líneas de pobreza del Banco Mundial y sus cifras, que se calculan con esas encuestas, son en buena parte su herencia.
  \item \emph{La nutrición}: en la India, el consumo de calorías por habitante bajó durante décadas de crecimiento; no por más pobreza, sino porque el trabajo pesado y las enfermedades disminuyeron y se diversificó la dieta \parencite{DeatonDreze2009}. Medir bien obliga a desconfiar de los indicadores sencillos.
  \item \emph{El escéptico}: el mismo \autor{Deaton} será uno de los críticos más duros de los experimentos (apartado~\ref{sec:valoracion-rct}).
  \end{itemize}
\item \emph{Nexo}: con \autor{Sen} y \autor{Deaton}, el foco pasa del país a las personas. La pregunta siguiente es cómo deciden los pobres y en qué mercados.
\end{itemize}

% --------------------------------------------------------------------------
\subsection{Comportamiento y mercados de los pobres}\label{sec:micro}

\subsubsection{Enfoque cultural: el papel de las aspiraciones y del comportamiento}\label{sec:cultural}

\begin{itemize}
\item \emph{Conclusión anticipada}: la pobreza no solo limita los recursos; también cambia cómo se decide, y eso puede perpetuarla.
\item \emph{Las aspiraciones}: para \textcite{Ray2006}, cada persona fija sus aspiraciones mirando a quienes tiene cerca. La \concepto{brecha de aspiraciones} es la distancia entre el nivel de vida al que se aspira y el actual:
  \begin{itemize}
  \item si es pequeña, no hay incentivo a esforzarse;
  \item si es enorme, tampoco, porque la meta parece inalcanzable;
  \item solo una brecha intermedia mueve a invertir.
  \end{itemize}
  \textcite{GenicotRay2017} muestran que así la desigualdad se reproduce: en una sociedad muy desigual, los pobres miran a metas tan lejanas que no invierten.
\item \emph{La desesperanza}: \textcite{Duflo2012Tanner} sostiene que la pobreza quita la confianza en poder cambiar las cosas, y que esa falta de esperanza frena la inversión en educación o en un negocio; parte del éxito de los programas de graduación (apartado~\ref{sec:evidencia}) se debería a que la devuelven.
\item \emph{La escasez}: preocuparse por llegar a fin de mes ocupa la atención y empeora las decisiones. Los mismos cultivadores de caña de azúcar de la India resolvían peor las pruebas de razonamiento antes de la cosecha, cuando eran pobres, que después \parencite{ManiEtAl2013}: no es que los pobres decidan peor, es que la pobreza agota \parencite{MullainathanShafir2013}.
\item \emph{Las trampas de los hogares}: la trampa de pobreza puede darse en una familia aunque no se dé en un país. Si la renta de mañana depende de la de hoy con forma de S, hay un umbral por debajo del cual el hogar se empobrece y por encima del cual prospera \parencite[cap.~1]{BanerjeeDuflo2011}: por la nutrición (quien come poco rinde poco y gana poco; \cite{DasguptaRay1986}), por los activos (sin una vaca o un taller solo hay trabajos mal pagados) o por el crédito.
  \begin{itemize}
  \item \emph{A favor}: en el programa de graduación de BRAC en Bangladés, que entregó activos (sobre todo vacas) a hogares muy pobres elegidos al azar, los que recibieron lo bastante para superar un umbral acumularon, cambiaron de ocupación y salieron de la pobreza durante 11 años; los que quedaron por debajo volvieron atrás \parencite{BalboniEtAl2022}.
  \item \emph{En contra, a escala de países}: casi todos los pobres crecen a largo plazo y los mecanismos son más continuos de lo que supone la teoría \parencite{KraayMcKenzie2014Do}. Las trampas existen para algunos hogares y lugares, pero no son la norma.
  \end{itemize}
\item \emph{La implicación}: las políticas tienen que contar con el comportamiento (simplificar, recordar, dar incentivos pequeños e inmediatos) \vertema{3.A.1}, y donde hay umbral, una ayuda grande y temporal vale más que una pequeña y permanente.
\item \emph{Nexo}: las aspiraciones y la escasez explican decisiones individuales; las toman, además, en mercados que fallan.
\end{itemize}

\subsubsection{Mercados que fallan: el crédito y los seguros}\label{sec:mercados}

\begin{itemize}
\item \emph{El crédito}: los pobres no tienen colateral, prestar en pequeño es caro y el prestamista no sabe quién es fiable ni puede vigilar lo que se hace con el dinero. El resultado es el racionamiento del crédito \parencite{StiglitzWeiss1981Credit} \vertema{3.A.13} y unos prestamistas informales con tipos de interés muy altos.
\item \emph{Los microcréditos con responsabilidad conjunta}: el Banco Grameen de Bangladés, fundado por \autor{Muhammad Yunus} (Premio Nobel de la Paz de 2006), presta a grupos de vecinas que responden unas por otras. Aprovecha lo que el banco no sabe y los vecinos sí:
  \begin{itemize}
  \item \emph{la selección}: los prestatarios seguros se agrupan entre sí y los arriesgados quedan juntos, lo que alivia la selección adversa \parencite{Ghatak1999Group,GhatakGuinnane1999};
  \item \emph{la vigilancia}: los miembros del grupo se vigilan unos a otros, lo que reduce el riesgo moral \parencite{Stiglitz1990Peer};
  \item \emph{las sanciones sociales}: quien no paga pierde la reputación y el acceso futuro al crédito.
  \end{itemize}
\item \emph{Los seguros}: casi no hay seguros formales (información asimétrica, poca seguridad jurídica, costes altos), y los hogares se aseguran con redes de parientes y vecinos, que cubren los choques de cada familia pero no los que golpean a todo el pueblo, como una sequía. Sin seguro, se evitan inversiones rentables pero arriesgadas\footnote{Por la misma falta de seguridad jurídica, el Grupo Banco Mundial tiene un organismo, el Organismo Multilateral de Garantía de Inversiones (OMGI, MIGA en inglés), que asegura a los inversores extranjeros contra los riesgos políticos en los países en desarrollo \vertema{3.B.30}.}.
\item \emph{La tecnología puede suplir al mercado}: en Kenia, el dinero móvil M-Pesa permitió enviar dinero a la familia ante un choque y ahorrar; su difusión sacó de la pobreza a unos 194.000 hogares, el 2~\% del total \parencite{SuriJack2016} (una estimación cuasiexperimental, no un experimento).
\item \emph{Nexo}: para saber qué programas corrigen estos fallos hay que medirlo, y la forma más creíble es el experimento.
\end{itemize}

% --------------------------------------------------------------------------
\subsection{La aproximación experimental}\label{sec:experimental}

\subsubsection{El método}\label{sec:metodo}


\begin{itemize}
\item \emph{La idea}: para saber si un programa funciona hay que comparar lo que les pasa a quienes lo reciben con lo que les habría pasado sin él, que nunca se observa. \autor{Kremer} en Kenia desde mediados de los 90, y \autor{Banerjee} y \autor{Duflo} en la India, llevaron a la economía del desarrollo el método de los ensayos clínicos: el \concepto{experimento aleatorizado controlado} (\emph{randomized controlled trial}, RCT). «Ver para creer»: antes de recomendar un programa, se prueba.
\item \emph{Los resultados potenciales} \parencite{Rubin1974}: queremos el efecto de un programa (unas pastillas, un libro de texto) sobre un resultado (la asistencia a la escuela). Cada persona $i$ tiene 2 resultados posibles, $Y_i(1)$ si lo recibe y $Y_i(0)$ si no ($D_i=1$ o $D_i=0$), pero solo se ve uno. Lo que interesa es el \concepto{efecto medio del tratamiento}, $\mathrm{ATE}\equiv\mathbb{E}[Y_i(1)-Y_i(0)]$.
\item \emph{El problema}: comparar a quienes reciben el programa con quienes no mezcla el efecto con las diferencias previas entre ellos:
\funcionclave[eq:sesgo-seleccion]{\underbrace{\mathbb{E}[Y_i|D_i=1]-\mathbb{E}[Y_i|D_i=0]}_{\text{lo que se observa}}=\underbrace{\mathbb{E}[Y_i(1)-Y_i(0)|D_i=1]}_{\text{efecto en los tratados}}+\underbrace{\mathbb{E}[Y_i(0)|D_i=1]-\mathbb{E}[Y_i(0)|D_i=0]}_{\text{sesgo de selección}}}
  Si los colegios que piden libros de texto son los de padres más implicados, sus alumnos habrían sacado mejores notas también sin libros: la comparación exagera el efecto. Es el \concepto{sesgo de selección}.
\item \emph{La solución}: si el programa se asigna \emph{al azar}, tratados y controles son iguales en promedio en todo (lo que se ve y lo que no) salvo en el programa; el sesgo de selección es nulo y la diferencia de medias estima el efecto medio \parencite{AngristPischke2009} (la demostración, en el anexo~\ref{anexo:evaluacion}).
\item \emph{Validez interna y externa}:
  \begin{itemize}
  \item \emph{interna}: la aleatorización elimina el sesgo de selección \emph{entre} tratados y controles: el efecto estimado es el de ese programa en esa muestra;
  \item \emph{externa}: otra cosa es si el resultado vale fuera de la muestra, y ahí el azar no ayuda, porque la muestra misma no es aleatoria. Si la población son todos los colegios de Kenia y la muestra, 100 colegios, el país y los colegios se eligen a menudo por conveniencia (cerca de la capital, donde hay un socio dispuesto), y solo participa quien quiere. Los sitios que aceptan ser evaluados suelen ser aquellos en los que el programa funciona mejor: en 111 evaluaciones de un programa de ahorro de energía en Estados Unidos, las primeras se hicieron donde mejor funcionaba y sobrestimaban su efecto en el resto \parencite{Allcott2015}.
  \end{itemize}
\item \emph{Otras dificultades}: no todos los asignados cumplen, y se estima entonces el efecto de la \emph{intención de tratar}; el programa puede afectar a los no tratados (contagios, precios); y lo que funciona a pequeña escala puede no hacerlo a gran escala \parencite{DufloGlennersterKremer2007}.
\item \emph{Cuando no se puede sortear}: los \emph{cuasiexperimentos} buscan una asignación casi aleatoria en los datos, como \textcite{Duflo2001} con la construcción de escuelas en Indonesia; el anexo~\ref{anexo:evaluacion} da la idea de cada técnica \vertema{4.B.6}.
\end{itemize}

\subsubsection{La evidencia experimental}\label{sec:evidencia}

\begin{itemize}
\item \emph{El volumen}: los investigadores de J-PAL, el laboratorio que fundaron \autor{Banerjee}, \autor{Duflo} y \autor{Sendhil Mullainathan} en 2003, han hecho más de 2.700 evaluaciones aleatorias en más de 140 países \parencite{JPAL2026}. Algunos resultados son ya clásicos (el anexo~\ref{anexo:experimentos} los ordena por áreas).
\item \emph{Salud: los precios importan muchísimo}. Muchos de los experimentos clásicos se hicieron en Kenia:
  \begin{itemize}
  \item la desparasitación en las escuelas de Busia redujo el absentismo en una cuarta parte, también en los colegios vecinos por el menor contagio \parencite{MiguelKremer2004}, y 20 años después quienes la recibieron consumían más y ganaban más por hora \parencite{HamoryEtAl2021};
  \item cobrar una tasa pequeña por los antiparasitarios redujo la participación en torno a un 80~\% respecto a darlos gratis \parencite{KremerMiguel2007}, y lo mismo pasa con las mosquiteras contra la malaria \parencite{CohenDupas2010}: cobrar no hace que se valore más, solo reparte menos.
  \item En la India, unos campamentos de vacunación fiables triplicaron la inmunización completa, del 6~\% al 18~\%, y regalar un kilo de lentejas por vacuna la llevó al 39~\%, a la mitad de coste por niño vacunado \parencite{BanerjeeDufloGlennersterKothari2010}.
  \end{itemize}
\item \emph{Educación: más recursos, por sí solos, ayudan poco}:
  \begin{itemize}
  \item en Kenia, los libros de texto solo mejoraron las notas de los alumnos que ya iban bien \parencite{GlewweKremerMoulin2009}, y las láminas para el aula no tuvieron efecto \parencite{GlewweKremerMoulinZitzewitz2004};
  \item las comidas gratis en los preescolares de Kenia aumentaron la asistencia un 31~\%, pero las pruebas solo mejoraron donde el maestro tenía experiencia \parencite{VermeerschKremer2005};
  \item lo que funciona es enseñar según el nivel real de los alumnos, como los refuerzos de la ONG Pratham en la India \parencite{BanerjeeColeDufloLinden2007}, que hoy llegan a millones de niños \parencite{BanerjeeEtAl2017}, y los incentivos del profesor: en la India, pagar según la asistencia comprobada con una cámara redujo a la mitad el absentismo de los maestros \parencite{DufloHannaRyan2012}.
  \end{itemize}
\item \emph{Crédito}: 6 experimentos en Bosnia y Herzegovina, Etiopía, la India, Marruecos, México y Mongolia muestran que los microcréditos tienen efectos modestos (algo más de inversión en negocios que ya existían, ninguna transformación del consumo ni de la renta) y una demanda menor de la esperada \parencite{BanerjeeKarlanZinman2015}. Y en Filipinas, quitar la responsabilidad conjunta no aumentó los impagos \parencite{GineKarlan2014}: las sanciones sociales pesan menos de lo que se creía.
\item \emph{Agricultura}: los agricultores del oeste de Kenia no compran fertilizante aunque les sale muy rentable, por \emph{sesgo al presente}: piensan comprarlo, pero cuando llega la siembra ya no tienen dinero. Ofrecérselo con un pequeño descuento justo tras la cosecha, cuando tienen liquidez, aumentó su uso tanto como subvenciones mucho mayores \parencite{DufloKremerRobinson2011}.
\item \emph{Transferencias}:
  \begin{itemize}
  \item las \emph{condicionadas} a la escolarización y a las revisiones médicas, como \emph{Progresa} en México (1997), que se extendió a los pueblos por sorteo, aumentaron la escolarización y mejoraron la salud infantil \parencite{Schultz2004}, y se copiaron en casi toda América Latina;
  \item las \emph{incondicionales} de GiveDirectly en Kenia aumentaron el consumo, los activos y el bienestar, sin gastarse en alcohol ni tabaco \parencite{HaushoferShapiro2016}; y cada dólar transferido generó unos 2,5 dólares de actividad en la comarca, con muy poca inflación \parencite{EggerEtAl2022};
  \item los \emph{programas de graduación}, que combinan un activo productivo, formación, un pequeño ingreso y acompañamiento, mejoraron el consumo, los activos y la salud mental de los más pobres en 6 países, años después de terminar \parencite{BanerjeeEtAl2015Graduacion}.
  \end{itemize}
\item \emph{El reconocimiento}: por estos trabajos recibieron \autor{Banerjee}, \autor{Duflo} y \autor{Kremer} el Premio Nobel de 2019 \parencite{NobelPrize2019}.
\end{itemize}

\subsubsection{Valoración}\label{sec:valoracion-rct}

\begin{itemize}
\item \emph{Ventajas}:
  \begin{itemize}
  \item \emph{credibilidad}: bien hecho, un experimento da una estimación sin sesgo de selección y sin supuestos difíciles de comprobar;
  \item \emph{política basada en la evidencia}: permite priorizar los programas que funcionan y comparar su coste por resultado;
  \item \emph{diseño antes que evaluación}: el investigador participa desde el principio y prueba variantes (precios, recordatorios, incentivos), con lo que aprende por qué funciona un programa y no solo si funciona; la evaluación tradicional llegaba después, con los datos que hubiera.
  \end{itemize}
\item \emph{Críticas}:
  \begin{itemize}
  \item \emph{la validez externa}: un resultado de Kenia puede no valer en Guatemala; sin una teoría que explique \emph{por qué} funciona un programa, no se sabe a qué otros sitios trasladarlo \parencite{Deaton2010,DeatonCartwright2018}. Los laureados responden repitiendo los experimentos en varios países y ligándolos a la teoría \parencite{BanerjeeDuflo2009}. «Del dicho al hecho hay mucho trecho»: del experimento a la política a escala;
  \item \emph{los mecanismos}: un experimento estima un efecto medio, pero no su reparto ni sus mecanismos, y el propio sorteo puede cambiar el comportamiento de quien participa \parencite{HeckmanSmith1995}; los modelos estructurales, que van de la teoría a los datos, sí permiten simular políticas nuevas;
  \item \emph{la jerarquía de la evidencia}: el experimento no es siempre lo mejor; con el mismo presupuesto, un estudio observacional bien hecho puede acercarse más a lo que se quiere saber que un experimento caro y pequeño \parencite{Ravallion2020};
  \item \emph{pensar en pequeño}: para \textcite{Deaton2010}, mejorar un programa sirve de poco si la economía está mal gobernada, con la macroeconomía y las instituciones en contra: es arreglar una ventana en una casa que se cae. \textcite{Rodrik2009} pide combinar los 2 enfoques, experimentar con un diagnóstico de qué restricción importa, en el debate entre pensar en grande y en pequeño \parencite{CohenEasterly2009};
  \item \emph{los efectos del propio experimento}: quien se sabe observado cambia su conducta (\emph{efecto Hawthorne}, aunque su versión original no se sostiene en los datos; \cite{LevittList2011});
  \item \emph{la ética}: se experimenta con personas muy pobres, y quien queda en el grupo de control se queda sin un programa que puede ser bueno \parencite{Deaton2020}.
  \end{itemize}
\item \emph{Balance}: los experimentos no explican el crecimiento de las naciones, pero sí qué programas concretos funcionan, y han llevado a la economía del desarrollo una exigencia de evidencia que ya no tiene vuelta atrás.
\end{itemize}

\begin{anotaciones}
\textbf{Para leer más.} El debate que abrió el Premio Nobel de 2019, en «This year's Nobel prizes prompt soul-searching among economists», \emph{The Economist}, 21 de noviembre de 2019: \url{https://www.economist.com/finance-and-economics/2019/11/21/this-years-nobel-prizes-prompt-soul-searching-among-economists} \parencite{Economist2019Nobel}.
\end{anotaciones}

% --------------------------------------------------------------------------
\subsection{Implicaciones para el diseño de políticas}\label{sec:politicas}

\begin{itemize}
\item \emph{Conclusión anticipada}: no hay una receta única, pero la evidencia acumulada deja 6 lecciones. «Obras son amores, que no buenas razones»: la política de desarrollo se juzga por lo que consigue, no por lo que promete.
\item \emph{1. Primero, el diagnóstico}: ni el gran empujón ni el Consenso de Washington sirvieron para todos; hay que buscar la restricción que más limita a cada país y empezar por ella \parencite{HausmannRodrikVelasco2008}.
\item \emph{2. La macro y las instituciones son condición necesaria}: estabilidad macroeconómica (lo que quedó del Consenso de Washington), propiedad segura, límites al poder y un Estado capaz de recaudar y prestar servicios, con instituciones adaptadas al país y a su distancia a la frontera (apartados~\ref{sec:ortodoxa}, \ref{sec:formal} y~\ref{sec:institucional}). Sin ellas, ningún programa micro compensa.
\item \emph{3. Evaluar antes de escalar}: las políticas se diseñan con su evaluación incorporada, con experimentos cuando se pueda, y lo que funciona se lleva a escala con cuidado, porque lo que funciona con una ONG puede fallar con la administración \parencite{BanerjeeEtAl2017}. Así trabajan J-PAL, Innovations for Poverty Action y el Banco Mundial, y la cooperación española \vertema{4.A.26} \vertema{4.B.6}.
\item \emph{4. Diseñar para personas reales}: precios cero en la prevención, incentivos pequeños e inmediatos (las lentejas), compromisos que eviten el sesgo al presente (el fertilizante tras la cosecha) y trámites sencillos (apartado~\ref{sec:cultural}).
\item \emph{5. Dar dinero también es una política, y a veces hace falta un empujón grande}: las transferencias, condicionadas o no, reducen la pobreza sin los efectos negativos que se temían, y donde hay umbrales los programas de graduación sacan de la pobreza a hogares que una ayuda pequeña dejaría donde estaban \parencite{BalboniEtAl2022}. La ayuda internacional funciona mejor ligada a la evidencia y a las instituciones que a la brecha de financiación \vertema{3.B.29} \vertema{3.B.30} \vertema{3.B.42}.
\item \emph{6. Una política industrial moderna}: centrada en fallos concretos (coordinación, autodescubrimiento), con apoyos condicionados a resultados y retirada de los que no funcionan \parencite{JuhaszLaneRodrik2024}; y, desde la renta media, pasar de invertir a incorporar tecnología y a innovar \parencite{BM2024Informe}.
\item \emph{El marco común}: la Agenda 2030 de las Naciones Unidas, con su objetivo 1, «poner fin a la pobreza en todas sus formas en todo el mundo» \parencite{ONU2015Agenda}, y la financiación del clima, porque los países pobres y cálidos son los que más pierden con el calentamiento, con responsabilidades comunes pero diferenciadas \vertema{3.B.31}.
\item \emph{Recapitulación}: la aproximación macroeconómica explica por qué el atraso de un país puede persistir (complementariedades, instituciones), y la microeconómica, por qué persiste el de un hogar (comportamiento, mercados que fallan) y qué programas funcionan. La política de desarrollo necesita las 2 cosas: pensar en grande para el marco y en pequeño para el diseño.
\end{itemize}


% ==========================================================================
\section*{Conclusión}
\minutos{1}

\begin{itemize}
\item \textbf{\emph{Recapitulación}}:
  \begin{itemize}
  \item En conclusión,
  \begin{itemize}
  \item los países en desarrollo comparten dualidad, mercados incompletos, poca capacidad del Estado y una transición demográfica a medias; el desarrollo se mide con el desarrollo humano, la renta y la pobreza, porque ninguna cifra lo resume;
  \item a escala de los países, la pobreza extrema ha caído del 43~\% al 10~\% desde 1990, pero se concentra en África. La teoría del desarrollo nació separada de la del crecimiento porque esta suponía convergencia y un solo sector, y acertó en 3 ideas: la dinámica económica depende de muchos factores, esos factores se refuerzan entre sí y hay margen para políticas activas, aunque son difíciles de identificar y de aplicar. Sus recetas generales, primero intervencionistas y después liberales, fallaron; las teorías recientes formalizaron sus intuiciones y señalan a las instituciones como la causa última más robusta;
  \item a escala de las personas, de \autor{Sen} y \autor{Deaton} a los experimentos, sabemos cómo deciden los pobres y qué programas funcionan, en un debate abierto entre el método estructural, que va de la teoría a los datos, y el experimental, que va de los datos a la teoría.
  \end{itemize}
  \end{itemize}
\item \textbf{\emph{Relevancia}}:
  \begin{itemize}
  \item Una de cada 10 personas del mundo, unos 800 millones, sigue en la pobreza extrema, y 108 países están en la renta media \parencite{BM2026PIP,BM2024Informe}: saber qué funciona decide cientos de millones de vidas.
  \end{itemize}
\item \textbf{\emph{Extensiones}}:
  \begin{itemize}
  \item Los modelos de crecimiento que formalizan estas ideas \vertema{3.A.43} \vertema{3.A.44} y su evidencia \vertema{3.A.45}; el comercio y el crecimiento \vertema{3.B.9}; la financiación exterior, la deuda y la ayuda \vertema{3.B.29}, los bancos de desarrollo \vertema{3.B.30}, la cooperación de la Unión Europea y de España \vertema{3.B.42} \vertema{4.A.26}, el clima \vertema{3.B.31}, la evaluación de las políticas públicas \vertema{4.B.6} y el método de la economía \vertema{3.A.1}.
  \end{itemize}
\item \textbf{\emph{Opinión}}:
  \begin{itemize}
  \item Los 2 enfoques premiados son complementarios.

  \item Ni China ni Corea del Sur se desarrollaron con experimentos, sino con mucho ahorro, inversión, exportaciones y un Estado que coordinaba\footnote{El ahorro nacional bruto de China llegó al 51~\% del PIB en 2008 y aún era del 42,9~\% en 2024 \parencite{FMI2026WEO}. China pasó de producir con mucho trabajo y poco capital a ser una potencia industrial, con un Estado que coordinó la inversión y amplió la capacidad productiva.}; pero dentro de ese marco, saber cómo responden las personas multiplica la eficacia de cada euro de ayuda.
  \end{itemize}
\item \textbf{\emph{Idea final}}:
  \begin{itemize}
  \item Volvamos al refrán: «unos nacen con estrella y otros estrellados».

  \item Nacer en Kenia o en España sigue siendo, en buena medida, cuestión de estrella.

  \item Pero la economía del desarrollo enseña que la estrella no es un destino: Botsuana, China o Corea del Sur lo muestran, y con instituciones que funcionen, una macroeconomía estable y políticas que se ha comprobado que funcionan, la estrella de nacimiento puede pesar cada vez menos.
  \end{itemize}
\end{itemize}

\clearpage
\appendix

% ==========================================================================
\section{Anexo: metodología del índice de desarrollo humano}\label{anexo:idh}

\begin{itemize}
\item \emph{Paso 1: los índices de cada dimensión}, normalizados entre un mínimo y un máximo fijos \parencite{PNUD2025NotasTecnicas}:
  \[
  I_{\mathrm{S}}=\frac{\text{esperanza de vida}-20}{85-20},\qquad
  I_{\mathrm{E}}=\frac{1}{2}\left(\frac{\text{años medios}}{15}+\frac{\text{años esperados}}{18}\right),
  \]
  \[
  I_{\mathrm{R}}=\frac{\ln(\mathrm{RNB}/L)-\ln 100}{\ln 75.000-\ln 100},
  \]
  con la RNB por habitante en dólares de PPA de 2021. Los topes tienen una justificación: ninguna sociedad del siglo XX tuvo una esperanza de vida inferior a 20 años; 18 años esperados equivalen a terminar un máster, y por encima de 75.000 dólares más renta apenas añade desarrollo humano. La renta va en logaritmos porque cada dólar adicional amplía menos las capacidades.
\item \emph{Paso 2: la agregación}. El IDH es la \emph{media geométrica} de los 3 índices,
  \[
  \mathrm{IDH}=\left(I_{\mathrm{S}}\cdot I_{\mathrm{E}}\cdot I_{\mathrm{R}}\right)^{1/3}.
  \]
  Es geométrica desde 2010 \parencite{PNUD2010IDH}; antes era aritmética. Así, un índice muy bajo pesa más y una dimensión no se compensa del todo con las otras. Como cada índice está entre 0 y 1, el IDH también. ($I_{\mathrm{S}}$, $I_{\mathrm{E}}$ e $I_{\mathrm{R}}$ son letras locales de este anexo.)
\item \emph{Grupos}: desarrollo humano muy alto (IDH de 0,800 o más), alto (de 0,700 a 0,799), medio (de 0,550 a 0,699) y bajo (menos de 0,550).
\item \emph{Datos de 2023}: lo encabeza Islandia (0,972), seguida de Suiza y Noruega (0,970); España ocupa el puesto 28 (0,918); China, el 78 (0,797); Kenia, el 143 (0,628), y lo cierra Sudán del Sur (0,388). La media mundial es 0,756 \parencite{PNUD2025IDH}. La evolución por grupos y regiones desde 1990 está en el gráfico~\ref{fig:idh-regiones}.
\item \emph{Otros índices del PNUD}: el IDH ajustado por la desigualdad, los índices de desarrollo de género y de desigualdad de género y el índice de pobreza multidimensional (con la Universidad de Oxford), que en 2010 sustituyó a los antiguos índices de pobreza humana, uno para los países en desarrollo y otro para los de la OCDE \parencite{PNUD2010IDH}.
\end{itemize}

% ==========================================================================
\section{Anexo: teorías clásicas del desarrollo en detalle}\label{anexo:clasicas}

\begin{itemize}
\item Este anexo recoge el detalle de las teorías del apartado~\ref{sec:clasicas} que no se canta: el modelo de crecimiento acumulativo de \autor{Kaldor}, los argumentos de la CEPAL, la comparación entre la sustitución de importaciones y la orientación a la exportación (que se estudia aparte con su evidencia \vertema{3.B.9}), los 10 puntos del Consenso de Washington y el gran empujón de \autor{Murphy}, \autor{Shleifer} y \autor{Vishny}.
\end{itemize}

\paragraph{El crecimiento acumulativo de \autor{Kaldor}}

\begin{itemize}
\item En la versión de \textcite{DixonThirlwall1975}, con $a_V$, $b_V$, $a_D$ y $b_D$ constantes positivas (letras locales de este anexo):
  \begin{align}
  \gamma_{Y/L}&=a_V+b_V\,\gamma_Y &&\text{(ley de \autor{Kaldor}-\autor{Verdoorn})},\label{eq:verdoorn}\\
  \gamma_Y&=a_D+b_D\,\gamma_{Y/L} &&\text{(la productividad abarata las exportaciones y tira de la demanda)}.\label{eq:demanda-kaldor}
  \end{align}
\item Sustituyendo \eqref{eq:verdoorn} en \eqref{eq:demanda-kaldor}, la tasa de crecimiento de equilibrio es
  \[
  \gamma_Y^*=\frac{a_D+b_D\,a_V}{1-b_D\,b_V}.
  \]
\item \emph{Si $b_D\,b_V<1$}, el equilibrio es estable: una perturbación se corrige y la economía vuelve a $\gamma_Y^*$. Si además las constantes son nulas ($a_V=a_D=0$), $\gamma_Y^*=0$: la realimentación no basta para crecer de forma sostenida, como en el modelo de \autor{Solow} sin progreso técnico, que acaba en un estado estacionario.
\item \emph{Si $b_D\,b_V>1$}, el equilibrio es inestable y la economía se aleja de él: crece cada vez más deprisa o se hunde. Las diferencias de partida se amplifican, que es la causación acumulativa de \autor{Myrdal}. Un impulso que eleve las constantes (la demanda exterior, la política industrial) puede poner a un país en la buena senda.
\end{itemize}

\paragraph{La CEPAL: por qué se rompió la correa de transmisión}

\begin{itemize}
\item En el siglo XIX, el centro transmitía su crecimiento a la periferia comprándole materias primas y alimentos y enviándole capital y trabajadores. En el XX dejó de hacerlo, por 3 razones:
  \begin{itemize}
  \item cayó la demanda relativa de productos primarios, por su baja elasticidad renta y por el proteccionismo agrario de Europa y Estados Unidos;
  \item los ingresos de exportación eran bajos y volátiles, con márgenes escasos;
  \item creció el comercio de servicios, que apenas usa materias primas.
  \end{itemize}
\item \emph{La industria naciente}: 2 fallos de mercado justifican protegerla \vertema{3.B.8}:
  \begin{itemize}
  \item \emph{un mercado de capitales imperfecto}: si el ahorro del sector tradicional no puede financiar los sectores nuevos, estos solo crecen con sus propios beneficios;
  \item \emph{la apropiabilidad}: los pioneros pagan costes de aprendizaje de los que se benefician quienes llegan después, y nadie quiere ser el primero.
  \end{itemize}
  La política óptima es corregir el fallo en su origen (un mercado de capitales que funcione, una compensación a los pioneros); el arancel es solo un segundo óptimo \vertema{3.A.22}.
\end{itemize}

\paragraph{Sustitución de importaciones frente a orientación a la exportación}

\begin{itemize}
\item \emph{2 diferencias}:
  \begin{itemize}
  \item \emph{protección frente a apertura}: los defensores de la ISI razonaban con un modelo de 2 bienes y 2 factores: si el país es abundante en trabajo y exporta primarios, la alternativa es la industria intensiva en capital, en la que no tiene ventaja comparativa y que solo puede crecer protegida. Los países asiáticos encontraron un tercer bien, las manufacturas intensivas en trabajo, en las que sí tenían ventaja, y que podían desarrollar con libre comercio; subieron de gama cuando acumularon capital humano y tecnología;
  \item \emph{dirigismo frente a mercado}: la ISI se apoyaba en la demanda interna y en la planificación; la orientación a la exportación se sometía a la disciplina de los mercados exteriores.
  \end{itemize}
\item \emph{Las otras razones del éxito asiático}: el ahorro alto, la educación, unas cuentas públicas saneadas, un tipo de cambio estable y una inflación baja, la inversión extranjera y, en Corea del Sur y Taiwán, una reforma agraria previa que repartió la tierra \parencite{BM1993Milagro}.
\end{itemize}

\paragraph{Los 10 puntos del Consenso de Washington}

\begin{table}[H]
\centering
\caption{Las 10 recomendaciones del Consenso de Washington}\label{tab:consenso}
\small
\renewcommand{\arraystretch}{1.25}
\begin{tabular}{>{\raggedright\arraybackslash}m{3.6cm}>{\raggedright\arraybackslash}m{11.6cm}}
\toprule
\rowcolor{tceeazulpalido}\textbf{Bloque} & \textbf{Recomendaciones}\\
\midrule
Estabilidad macroeconómica & 1) disciplina fiscal; 2) reorientar el gasto público hacia la sanidad, la educación y las infraestructuras; 3) reforma tributaria con bases amplias y tipos marginales moderados\\
Precios de mercado & 4) tipos de interés fijados por el mercado; 5) tipo de cambio competitivo\\
Apertura & 6) liberalización comercial; 7) apertura a la inversión directa extranjera\\
Mercado y Estado & 8) privatización de las empresas públicas; 9) desregulación (eliminar barreras de entrada y de salida); 10) derechos de propiedad seguros\\
\bottomrule
\end{tabular}

\smallskip
\fuente{elaboración propia a partir de \textcite{Williamson1990}.}
\end{table}

\paragraph{El gran empujón de \autor{Murphy}, \autor{Shleifer} y \autor{Vishny}}

\begin{itemize}
\item \emph{Supuestos} (letras locales de este anexo): $L$ trabajadores y $N$ sectores con el mismo peso en el gasto, de modo que cada bien recibe $1/N$ de la renta. En cada sector:
  \begin{itemize}
  \item la técnica \emph{tradicional} usa 1 unidad de trabajo por unidad de producto, en competencia, con salario 1 (el numerario) y precio 1;
  \item la \emph{moderna}, que solo puede usar una empresa por sector, necesita $F+c\,q$ unidades de trabajo para producir $q$, con $c<1$, y paga un salario $w>1$, la prima que hay que pagar para llevar a los trabajadores a la fábrica; vende al precio 1 para no perder el mercado frente a la tradicional.
  \end{itemize}
\item \emph{El beneficio de una empresa moderna} que vende $q$ es
  \[
  B(q)=q-w\,(F+c\,q)=(1-c\,w)\,q-w\,F ,\qquad\text{con } 1-c\,w>0 .
  \]
\item \emph{Si se moderniza sola}, la renta del país sigue siendo $L$ y vende $q^A=L/N$. \emph{Si se modernizan todas}, todo el trabajo está en las fábricas, $L=N(F+c\,q)$, y cada una vende $q^B=(L/N-F)/c$, que es mayor que $q^A$ si $(1-c)L/N>F$.
\item \emph{Hay equilibrios múltiples} si
  \[
  B(q^A)<0\leq B(q^B):
  \]
  ninguna empresa se moderniza sola, pero si lo hacen todas, todas ganan. El equilibrio tradicional es una trampa: el moderno es mejor para todos (más renta y beneficios no negativos), pero nadie da el primer paso. La prima salarial es la clave: con $w=1$, si a una empresa sola no le compensa modernizarse, tampoco les compensa a todas, y solo hay un equilibrio.
\item \emph{Gráficamente}, en el plano empleo-producción de un sector, paso a paso:
  \begin{itemize}
  \item \emph{las 2 técnicas y los 2 puntos} (gráfico~\ref{fig:gran-empujon-a}): la técnica tradicional es la recta $q=L_j$ y la moderna, la recta más inclinada $q=(L_j-F)/c$, que empieza en el coste fijo $F$. El punto $A$ es la producción de una empresa moderna sola ($q^A$) y el $B$, la de cada una si se modernizan todas ($q^B$);
  \item \emph{los salarios} (gráfico~\ref{fig:gran-empujon}): los costes de la empresa moderna se leen en las rectas de costes, $q=w\,L_j$, de pendiente $w$: sobre ellas el beneficio es nulo. Con $w_2$ la recta pasa entre $A$ y $B$: la empresa sola pierde y todas juntas ganan, y hace falta el gran empujón. Con un salario bajo ($w_1$) a la empresa le compensa modernizarse sola; con uno alto ($w_3$), ni siquiera todas juntas.
  \end{itemize}
\item \emph{La política}: una subvención temporal, un plan de inversión o, simplemente, unas expectativas coordinadas («todos se van a modernizar») llevan a la economía al buen equilibrio, y después se pueden retirar.
\end{itemize}

\begin{figure}[H]
\centering
\caption{El gran empujón (a): las 2 técnicas y los puntos $A$ y $B$}\label{fig:gran-empujon-a}
\grafico[0.55\textwidth]{gran-empujon-paso1}
\fuente{elaboración propia a partir de \textcite{MurphyShleiferVishny1989Industrialization} y \textcite[cap.~2]{Krugman1995Development}. Primer paso del gráfico~\ref{fig:gran-empujon}.}
\end{figure}

\begin{figure}[H]
\centering
\caption{El gran empujón (b): 2 equilibrios en el modelo de \autor{Murphy}, \autor{Shleifer} y \autor{Vishny}}\label{fig:gran-empujon}
\grafico[0.6\textwidth]{gran-empujon}
\fuente{elaboración propia a partir de \textcite{MurphyShleiferVishny1989Industrialization} y \textcite[cap.~2]{Krugman1995Development}.}
\end{figure}

% ==========================================================================
\section{Anexo: la evaluación de impacto paso a paso}\label{anexo:evaluacion}

\paragraph{Por qué la aleatorización elimina el sesgo de selección}

\begin{itemize}
\item \emph{Paso 1}: lo que se observa de cada persona es $Y_i=D_i\,Y_i(1)+(1-D_i)\,Y_i(0)$: el resultado con el programa si lo recibe y sin él si no. Por tanto, la diferencia de medias observada entre tratados y no tratados es
  \[
  \mathbb{E}[Y_i\mid D_i=1]-\mathbb{E}[Y_i\mid D_i=0]=\mathbb{E}[Y_i(1)\mid D_i=1]-\mathbb{E}[Y_i(0)\mid D_i=0].
  \]
\item \emph{Paso 2}: sumando y restando $\mathbb{E}[Y_i(0)\mid D_i=1]$, lo que los tratados habrían obtenido sin el programa (que no se observa):
  \begin{align*}
  \mathbb{E}[Y_i\mid D_i=1]-\mathbb{E}[Y_i\mid D_i=0]
  &=\mathbb{E}[Y_i(1)\mid D_i=1]-\mathbb{E}[Y_i(0)\mid D_i=1]\\
  &\quad+\mathbb{E}[Y_i(0)\mid D_i=1]-\mathbb{E}[Y_i(0)\mid D_i=0]\\
  &=\underbrace{\mathbb{E}[Y_i(1)-Y_i(0)\mid D_i=1]}_{\text{efecto en los tratados}}+\underbrace{\mathbb{E}[Y_i(0)\mid D_i=1]-\mathbb{E}[Y_i(0)\mid D_i=0]}_{\text{sesgo de selección}},
  \end{align*}
  que es la descomposición \eqref{eq:sesgo-seleccion}.
\item \emph{Paso 3}: si el programa se asigna al azar, $D_i$ es independiente de los resultados potenciales, de modo que $\mathbb{E}[Y_i(0)\mid D_i=1]=\mathbb{E}[Y_i(0)\mid D_i=0]=\mathbb{E}[Y_i(0)]$, y lo mismo con $Y_i(1)$. El sesgo de selección es cero y el efecto en los tratados coincide con el efecto medio:
  \[
  \mathbb{E}[Y_i\mid D_i=1]-\mathbb{E}[Y_i\mid D_i=0]=\mathbb{E}[Y_i(1)-Y_i(0)]=\mathrm{ATE}.
  \]
\item \emph{Paso 4}: se estima con la diferencia de las medias muestrales o, lo que es lo mismo, con la regresión por mínimos cuadrados $Y_i=b_0+b_1D_i+u_i$, donde $b_1$ estima el ATE ($b_0$, $b_1$ y $u_i$ son letras locales). Añadir variables de control no cambia lo que se estima, pero lo hace más preciso. \findemo
\item \emph{Lo que la demostración no dice}: la igualdad vale para la población de la que se sacó la muestra; de otras poblaciones no dice nada (la validez externa, apartado~\ref{sec:metodo}).
\end{itemize}

\paragraph{Cuando no todos cumplen: la intención de tratar}

\begin{itemize}
\item Si se sortea la \emph{oferta} del programa pero no todos los elegidos lo usan (ni todos los demás se quedan sin él), comparar a los asignados con los no asignados da el efecto de la \emph{intención de tratar}, que sigue sin sesgo pero mezcla el efecto con la participación. Dividiéndolo por la diferencia en la proporción de tratados entre los 2 grupos se obtiene el efecto en quienes participan porque se les ofrece, usando el sorteo como variable instrumental \parencite{AngristPischke2009}.
\end{itemize}

\paragraph{Cuando no se puede sortear: los cuasiexperimentos}

\begin{itemize}
\item \emph{Diferencias en diferencias}: se compara el cambio en el tiempo de los tratados con el de los no tratados, suponiendo que sin el programa habrían evolucionado en paralelo. \textcite{Duflo2001} lo aplica a la construcción de más de 61.000 escuelas en Indonesia en los 70, comparando las regiones con más y menos escuelas y a los niños que tenían edad de aprovecharlas con los mayores.
\item \emph{Regresión discontinua}: si el programa se da a quien está por debajo de un umbral (de renta, de nota), se compara a los que quedan justo a un lado y al otro, que son casi iguales.
\item \emph{Variables instrumentales}: una variable que afecta a quién recibe el programa pero no al resultado por otra vía (como la mortalidad de los colonos en el apartado~\ref{sec:institucional}).
\item \emph{Experimentos naturales}: un hecho ajeno a las personas (una lotería, una frontera, una reforma) que reparte el tratamiento casi al azar.
\item Su estudio detallado, con sus supuestos y la evaluación de las políticas públicas en España, se hace aparte \vertema{4.B.6}.
\end{itemize}

% ==========================================================================
\section{Anexo: experimentos de desarrollo por áreas}\label{anexo:experimentos}

\begin{table}[H]
\centering
\caption{Algunos experimentos clásicos de la economía del desarrollo}\label{tab:experimentos}
\footnotesize
\renewcommand{\arraystretch}{1.3}
\begin{tabular}{>{\raggedright\arraybackslash}m{1.9cm}>{\raggedright\arraybackslash}m{4.3cm}>{\raggedright\arraybackslash}m{5.4cm}>{\raggedright\arraybackslash}m{2.9cm}}
\toprule
\rowcolor{tceeazulpalido}\textbf{Área} & \textbf{Programa y lugar} & \textbf{Resultado} & \textbf{Fuente}\\
\midrule
Salud & Desparasitación en las escuelas (Busia, Kenia) & Una cuarta parte menos de absentismo, también en los colegios vecinos; más consumo y salario 20 años después & \textcite{MiguelKremer2004}, \textcite{HamoryEtAl2021}\\
 & Pago pequeño por los antiparasitarios (Kenia) & La participación cae en torno a un 80~\% & \textcite{KremerMiguel2007}\\
 & Mosquiteras gratis o a precio reducido (Kenia) & Cobrar reduce mucho la demanda y no mejora el uso & \textcite{CohenDupas2010}\\
 & Campamentos de vacunación y lentejas (Udaipur, India) & Inmunización completa: 6~\%, 18~\% y 39~\%, a mitad de coste & \textcite{BanerjeeDufloGlennersterKothari2010}\\
\midrule
Educación & Libros de texto y láminas (Kenia) & Sin efecto medio; los libros solo ayudan a los mejores & \textcite{GlewweKremerMoulin2009}, \textcite{GlewweKremerMoulinZitzewitz2004}\\
 & Comidas en los preescolares (Kenia) & +31~\% de asistencia; mejores pruebas solo con maestros con experiencia & \textcite{VermeerschKremer2005}\\
 & Refuerzo según el nivel (Pratham, India) & Fuerte mejora del aprendizaje de los rezagados; llevado a escala & \textcite{BanerjeeColeDufloLinden2007}, \textcite{BanerjeeEtAl2017}\\
 & Maestros contratados y grupos por nivel (Kenia) & Más aprendizaje que reducir alumnos por profesor fijo & \textcite{DufloDupasKremer2011,DufloDupasKremer2015}\\
 & Cámaras y pago según la asistencia (India) & El absentismo de los maestros se reduce a la mitad & \textcite{DufloHannaRyan2012}\\
\midrule
Crédito & Microcréditos (6 países) & Efectos modestos y poca demanda & \textcite{BanerjeeKarlanZinman2015}\\
 & Sin responsabilidad conjunta (Filipinas) & No aumentan los impagos & \textcite{GineKarlan2014}\\
\midrule
Agricultura & Fertilizante con descuento tras la cosecha (Kenia) & Tanto uso como con subvenciones mucho mayores: sesgo al presente & \textcite{DufloKremerRobinson2011}\\
\midrule
Transferencias & Progresa (México) & Más escolarización y mejor salud infantil & \textcite{Schultz2004}\\
 & GiveDirectly (Kenia) & Más consumo, activos y bienestar; multiplicador local de unos 2,5 & \textcite{HaushoferShapiro2016}, \textcite{EggerEtAl2022}\\
 & Programas de graduación (6 países; BRAC en Bangladés) & Efectos duraderos; umbral de activos & \textcite{BanerjeeEtAl2015Graduacion}, \textcite{BalboniEtAl2022}\\
\bottomrule
\end{tabular}

\smallskip
\fuente{elaboración propia a partir de los artículos citados.}
\end{table}

% ==========================================================================
\section{Anexo: preguntas de test}\label{anexo:test}

\begin{anotaciones}
Preguntas tipo test del primer ejercicio sobre la medición del desarrollo, las teorías clásicas del desarrollo, el Consenso de Washington, las instituciones y la aproximación experimental: \href{https://www.victorgutierrezmarcos.es/oposicion/temario/primer-ejercicio/test/simulador.html}{simulador de test} \parencite{GutierrezMarcos2026Simulador}.
\end{anotaciones}

\bibliografia
\end{document}
